\subsection{存在定理}

我们保留上一节的假设与记号。回顾，若 \(\alpha, \beta \in \mathrm{B}\) 且 \(\alpha \neq \beta\) ，则 \(n(\beta, \alpha) \leq 0\) ；而且，若 \(n(\beta, \alpha) = 0\) ，则 \(n(\alpha, \beta) = 0\) （第六章 §1 第 1 节 公式 (8)）。对 B 中两个不同元素组成的偶 \((\alpha, \beta)\) ，置

\[
x _ {\alpha \beta} = (\mathrm{ad} x _ {\alpha}) ^ {1 - n (\beta , \alpha)} x _ {\beta} y _ {\alpha \beta} = (\mathrm{ad} x _ {- \alpha}) ^ {1 - n (\beta , \alpha)} x _ {- \beta}.
\]

于是 \(x_{\alpha \beta} \in \mathfrak{a}_{+}, y_{\alpha \beta} \in \mathfrak{a}_{-}\) 。若用 \(\theta\) 表示 \(\mathfrak{a}\) 的典则自同构，则 \(\theta(x_{\alpha \beta}) = y_{\alpha \beta}\) 。

引理 6. 设 \(\alpha, \beta \in \mathrm{B}\) 且 \(\alpha \neq \beta\) 。则

\[
[ \mathfrak {a} _ {+}, y _ {\alpha \beta} ] = 0 \quad [ \mathfrak {a} _ {-}, x _ {\alpha \beta} ] = 0.
\]

第二个公式可由第一个公式利用自同构 \(\theta\) 得出。为证明第一个公式，只需证明 \([x_{\gamma}, y_{\alpha \beta}] = 0\) 对一切 \(\gamma \in B\) 成立。我们区分三种情形。

情形 1：\(\gamma \neq \alpha\) 且 \(\gamma \neq \beta\) 。在此情形，\(x_{\gamma}\) 与 \(x_{-\alpha}\) 及 \(x_{-\beta}\) 交换，从而与 \(y_{\alpha \beta}\) 交换。

情形 2：\(\gamma = \beta\) 。在此情形，\(x_{\gamma}\) 与 \(x_{-\alpha}\) 交换，故

\[
\begin{array}{r l} & {[ x _ {\gamma}, y _ {\alpha \beta} ] = (\mathrm{ad} x _ {- \alpha}) ^ {1 - n (\beta , \alpha)} [ x _ {\gamma}, x _ {- \beta} ]} \\ & {\qquad = - (\mathrm{ad} x _ {- \alpha}) ^ {1 - n (\beta , \alpha)} h _ {\beta} = - n (\alpha , \beta) (\mathrm{ad} x _ {- \alpha}) ^ {- n (\beta , \alpha)} x _ {- \alpha}.} \end{array}
\]

若 \(n(\beta, \alpha) < 0\) ，此表达式为零，因为 (ad \(x_{-\alpha}\) ). \(x_{-\alpha} = 0\) 。若 \(n(\beta, \alpha) = 0\) ，则 \(n(\alpha, \beta) = 0\) 。在两种情形下，都有 \([x_{\gamma}, y_{\alpha\beta}] = 0\) 。

情形 3：\(\gamma = \alpha\) 。在 \(\mathfrak{a}\) 的自同态代数中，

\[
[ - \operatorname{ad} h _ {\alpha}, \operatorname{ad} x _ {- \alpha} ] = 2 \operatorname{ad} x _ {- \alpha}
\]

且 {[}ad \(x_{\alpha}\) , ad \(x_{-\alpha}\) {]} = -ad \(h_{\alpha}\) ；于是由 §1 引理 1，

\[
[ \operatorname{ad} x _ {\alpha}, (\operatorname{ad} x _ {- \alpha}) ^ {1 - n (\beta , \alpha)} ] = (1 - n (\beta , \alpha)) (\operatorname{ad} x _ {- \alpha}) ^ {- n (\beta , \alpha)} (- \operatorname{ad} h _ {\alpha} - n (\beta , \alpha)).
\]

因此，

\[
\begin{array}{c} [ x _ {\gamma}, y _ {\alpha \beta} ] = [ \text {ad} x _ {\alpha}, (\text {ad} x _ {- \alpha}) ^ {1 - n (\beta , \alpha)} ] x _ {- \beta} + (\text {ad} x _ {- \alpha}) ^ {1 - n (\beta , \alpha)} (\text {ad} x _ {\alpha}) x _ {- \beta} \\ = - (1 - n (\beta , \alpha)) (\text {ad} x _ {- \alpha}) ^ {- n (\beta , \alpha)} (\text {ad} h _ {\alpha} + n (\beta , \alpha)) x _ {- \beta} \\ + (\text {ad} x _ {- \alpha}) ^ {1 - n (\beta , \alpha)} (\text {ad} x _ {\alpha}) x _ {- \beta}. \end{array}
\]

而 \([h_{\alpha},x_{-\beta}] + n(\beta ,\alpha)x_{-\beta} = 0\) 且 \([x_{\alpha},x_{-\beta}] = 0\) ，故 \([x_{\gamma},y_{\alpha \beta}] = 0\) 。

引理 7. \(\mathfrak{n}\) （\(\mathfrak{a}_{+}\) 中由诸 \(x_{\alpha \beta} (\alpha, \beta \in \mathrm{B}, \alpha \neq \beta)\) 生成的理想）是 \(\mathfrak{a}\) 的理想。\(\mathfrak{a}_{-}\) 中由诸 \(y_{\alpha \beta} (\alpha, \beta \in \mathrm{B}, \alpha \neq \beta)\) 生成的理想是 \(\mathfrak{a}\) 的理想，且等于 \(\theta(\mathfrak{n})\) 。

设 \(\mathfrak{n}' = \sum_{\alpha, \beta \in \mathrm{B}, \alpha \neq \beta} kx_{\alpha \beta}\) 。由于每个 \(x_{\alpha \beta}\) 是 \(\mathfrak{a}\) 中的齐次元素，\([\mathfrak{a}^0, \mathfrak{n}'] \subset \mathfrak{n}'\) （引理 5 与命题 2）。设 U （相应地 V ）是 \(\mathfrak{a}\) （相应地 \(\mathfrak{a}_+\) ）的包络代数，\(\sigma\) 是 U 在 \(\mathfrak{a}\) 上由 \(\mathfrak{a}\) 的伴随表示所定义的表示。\(\mathfrak{a}\) 的由 \(\mathfrak{n}'\) 生成的理想是 \(\sigma(\mathrm{U})\mathfrak{n}'\) 。而 \(\mathfrak{a} = \mathfrak{a}_+ + \mathfrak{a}^0 + \mathfrak{a}_-\) （命题 2），\(\sigma(\mathfrak{a}_-)\mathfrak{n}' = 0\) （引理 6），且由前所述 \(\sigma(\mathfrak{a}^0)\mathfrak{n}' \subset \mathfrak{n}'\) 。由 Poincaré-Birkhoff-Witt 定理，\(\sigma(\mathrm{U})\mathfrak{n}' = \sigma(\mathrm{V})\mathfrak{n}'\) ，这就证明了引理的第一个断言。由此推出，\(\theta(\mathfrak{a}_+) = \mathfrak{a}_-\) 中由诸 \(\theta(x_{\alpha \beta}) = y_{\alpha \beta}\) （ \(\alpha, \beta \in \mathrm{B}, \alpha \neq \beta\) ）生成的理想是理想 \(\theta(\mathfrak{n})\) （\(\mathfrak{a}\) 中的）。证毕。

a 的理想 \(\mathfrak{n} + \theta(\mathfrak{n})\) 是分次的，因为它由齐次元素生成。因此，李代数 \(\mathfrak{a}/(\mathfrak{n}+\theta(\mathfrak{n}))\) 是一个 Q-分次李代数；在本节其余部分，将它记作 \(g_{B}\) ，或简记作 g。由命题 2，若 \(g^{\mu} \neq 0\) ，则 \(\mu \in Q_{+}\) ，或 \(\mu \in Q_{-}\) ，或 \(\mu = 0\) 。用 \(X_{\alpha}\) （相应地 \(H_{\alpha}, X_{-\alpha}\) ）表示 \(x_{\alpha}\) （相应地 \(h_{\alpha}, x_{-\alpha}\) ）在 g 中的典则像。鉴于 a ，n 与 \(\theta(\mathfrak{n})\) 的定义，由此推出，g 是由生成族 \(((X_{\alpha}, H_{\alpha}, X_{-\alpha}))_{\alpha \in B}\) 和诸关系式所定义的李代数

\[
[ H _ {\alpha}, H _ {\beta} ] = 0\tag{16}
\]

\[
[ H _ {\alpha}, X _ {\beta} ] - n (\beta , \alpha) X _ {\beta} = 0\tag{17}
\]

\[
[ H _ {\alpha}, X _ {- \beta} ] + n (\beta , \alpha) X _ {- \beta} = 0\tag{18}
\]

\[
[ X _ {\alpha}, X _ {- \alpha} ] + H _ {\alpha} = 0\tag{19}
\]

\[
[ X _ {\alpha}, X _ {- \beta} ] = 0 (\alpha \neq \beta)\tag{20}
\]

\[
(\operatorname{ad} X _ {\alpha}) ^ {1 - n (\beta , \alpha)} X _ {\beta} = 0 \quad (\alpha \neq \beta)\tag{21}
\]

\[
(\operatorname{ad} X _ {- \alpha}) ^ {1 - n (\beta , \alpha)} X _ {- \beta} = 0 \quad (\alpha \neq \beta).\tag{22}
\]

设 \(z \in \mathfrak{g}\) ，\(\mu \in \mathrm{Q}\) 。则 \(z \in \mathfrak{g}^{\mu}\) 当且仅当 \([H_{\alpha}, z] = \langle \mu, \alpha^{\vee} \rangle z\) 对一切 \(\alpha \in \mathrm{B}\) 成立。这由引理 5 得出。

由于 \(\mathfrak{a}^0\cap (\mathfrak{n} + \theta (\mathfrak{n})) = 0\) ，从 \(\mathfrak{a}^0\) 到 \(\mathfrak{g}^0\) 的典则映射是同构。因此，\((H_{\alpha})_{\alpha \in \mathrm{B}}\) 是向量空间 \(\mathfrak{g}^0\) 的一个基。\(\mathfrak{g}^0\) （\(\mathfrak{g}_{\mathrm{B}}\) 的交换子代数）将记作 \(\mathfrak{h}_{\mathrm{B}}\) ，或简记作 \(\mathfrak{h}\) 。存在唯一的同构 \(\mu \mapsto \mu_{\mathrm{B}}\) （从 V 到 \(\mathfrak{h}^*\) ），使 \(\langle \mu_{\mathrm{B}},H_{\alpha}\rangle = \langle \mu ,\alpha^{\vee}\rangle\) 对一切 \(\mu \in \mathrm{V}\) 与一切 \(\alpha \in \mathrm{B}\) 成立。

a 的对合自同构 \(\theta\) 通过过渡到商定义了 g 的一个对合自同构，它同样记作 \(\theta\) 。我们有 \(\theta(X_{\alpha}) = X_{-\alpha}\) （对 \(\alpha \in B \cup (-B)\) ），以及 \(\theta(H_{\alpha}) = -H_{\alpha}\) 。

定理 1. 设 R 是一个既约根系，B 是 R 的一个基。设 g 是由生成族 \(((X_{\alpha}, H_{\alpha}, X_{-\alpha}))_{\alpha \in \mathrm{B}}\) 和关系式 (16) 至 (22) 所定义的李代数。设 \(h = \sum_{\alpha \in B} kH_{\alpha}\) 。则 \((\mathfrak{g}, \mathfrak{h})\) 是一个分裂半单李代数。同构 \(\mu \mapsto \mu_{B}\) （从 V 到 \(h^{*}\) ）把 R 映为 \((\mathfrak{g}, \mathfrak{h})\) 的根系。对一切 \(\mu \in R\) ，\(g^{\mu}\) 是相对于根 \(\mu\) 的特征子空间。

证明将循引理 8、9、10、11 的途径进行。

引理 8. 设 \(\alpha \in \mathrm{B} \cup (-\mathrm{B})\) 。则 \(\operatorname{ad} X_{\alpha}\) 是局部幂零的。 \footnote{一个自同态 \(u\) （向量空间 \(V\) 上的）称为局部幂零的（或殆幂零的），如果对每个 \(v \in V\) ，存在正整数 \(n\) 使 \(u^n(v) = 0\) （参见交换代数 第四章 §1 第 4 节 定义 2）。此时 \(\exp(u)\) 或 \(e^u\) 由公式 \(e^u(v) = \sum_{n \geq 0} (1/n!)u^n(v)\) 对一切 \(v \in V\) 定义。}

假设 \(\alpha \in \mathrm{B}\) 。设 \(\mathfrak{g}'\) 是这样的 \(z \in \mathfrak{g}\) 组成的集合：\((\operatorname{ad} X_{\alpha})^p z = 0\) 对充分大的 \(p\) 成立。由于 \(\operatorname{ad} X_{\alpha}\) 是 \(\mathfrak{g}\) 的导子，\(\mathfrak{g}'\) 是 \(\mathfrak{g}\) 的子代数。由 (21)，\(X_{\beta} \in \mathfrak{g}'\) 对一切 \(\beta \in \mathrm{B}\) 成立。由 (17)、(19)、(20)，\(H_{\beta} \in \mathfrak{g}'\) 与 \(X_{-\beta} \in \mathfrak{g}'\) 对一切 \(\beta \in \mathrm{B}\) 成立。于是 \(\mathfrak{g}' = \mathfrak{g}\) ，从而 \(\operatorname{ad} X_{\alpha}\) 是局部幂零的。由于 \(\operatorname{ad} X_{-\alpha} = \theta (\operatorname{ad} X_{\alpha})\theta^{-1}\) ，可见 \(\operatorname{ad} X_{-\alpha}\) 是局部幂零的。

我们将会看到 g 是有限维的，从而 \(ad X_{\alpha}\) 实际上是幂零的。

引理 9. 设 \(\mu, \nu \in \mathbf{Q}\) ，\(w \in \mathrm{W}(\mathrm{R})\) 满足 \(w\mu = \nu\) 。存在 \(\mathfrak{g}\) 的自同构把 \(\mathfrak{g}^{\mu}\) 变到 \(\mathfrak{g}^{\nu}\) 。

对一切 \(\alpha \in \mathrm{B}\) ，设 \(s_{\alpha}\) 是 V 的由 \(\alpha\) 定义的反射。由于 \(\mathrm{W}(\mathrm{R})\) 由诸 \(s_{\alpha}\) 生成（第六章 §1 第 5 节 注 1），只需就 \(w = s_{\alpha}\) 的情形证明引理。鉴于引理 8，我们可以定义

\[
\theta_ {\alpha} = e ^ {\mathrm{ad} X _ {\alpha}} e ^ {\mathrm{ad} X _ {- \alpha}} e ^ {\mathrm{ad} X _ {\alpha}}.
\]

与第一章 §6 第 8 节一样地验证，\(\theta_{\alpha}\) 是 \(\mathfrak{g}\) 的自同构。我们有

\[
\begin{array}{l} \theta_ {\alpha} (H _ {\alpha}) = e ^ {\mathrm{ad} X _ {\alpha}} e ^ {\mathrm{ad} X _ {- \alpha}} (H _ {\beta} - n (\alpha , \beta) X _ {\alpha}) \\ \qquad = e ^ {\mathrm{ad} X _ {\alpha}} \bigg (H _ {\beta} - n (\alpha , \beta) X _ {\alpha} + n (\alpha , \beta) X _ {- \alpha} - n (\alpha , \beta) H _ {\alpha} - \frac {n (\alpha , \beta)}{2} 2 X _ {- \alpha} \bigg) \\ \qquad = e ^ {\mathrm{ad} X _ {\alpha}} \left(H _ {\beta} - n (\alpha , \beta) H _ {\alpha} - n (\alpha , \beta) X _ {\alpha}\right) \\ \qquad = H _ {\beta} - n (\alpha , \beta) H _ {\alpha} - n (\alpha , \beta) X _ {\alpha} - n (\alpha , \beta) X _ {\alpha} - n (\alpha , \beta) (- 2 X _ {\alpha}) \\ \qquad = H _ {\beta} - n (\alpha , \beta) H _ {\alpha}. \end{array}
\]

若 \(z\in \mathfrak{g}^{\mu}\)

\[
\begin{array}{r l} & {[ H _ {\beta}, \theta_ {\alpha} ^ {- 1} z ] = \theta_ {\alpha} ^ {- 1} [ H _ {\beta} - n (\alpha , \beta) H _ {\alpha}, z ]} \\ & {\qquad = \theta_ {\alpha} ^ {- 1} (\langle \mu , \beta^ {\vee} \rangle z - n (\alpha , \beta) \langle \mu , \alpha^ {\vee} \rangle z)} \\ & {\qquad = \langle \mu - \langle \alpha^ {\vee}, \mu \rangle \alpha , \beta^ {\vee} \rangle \theta_ {\alpha} ^ {- 1} z = \langle s _ {\alpha} \mu , \beta^ {\vee} \rangle \theta_ {\alpha} ^ {- 1} z,} \end{array}
\]

故 \(\theta_{\alpha}^{-1}z\in\mathfrak{g}^{s_{\alpha}\mu}\) 。这表明 \(\theta_{\alpha}^{-1}\mathfrak{g}^{\mu}\subset\mathfrak{g}^{s_{\alpha}\mu}\) 。由于 \(\theta_{\alpha}\) 是自同构，且此包含关系对一切 \(\mu\in Q\) 成立，可见 \(\theta_{\alpha}^{-1}\mathfrak{g}^{\mu}=\mathfrak{g}^{s_{\alpha}\mu}\) ，这就证明了引理。

引理 10. 设 \(\mu\in Q\) ，并假设 \(\mu\) 不是根的倍数。存在 \(w\in\mathrm{W}(\mathrm{R})\) ，使 \(w\mu\) 关于基 B 的坐标中有一些 \(> 0\) ，另一些 \(< 0\)。

设 \(V_{R}\) 是向量空间 \(Q \otimes_{Z} R\) ，R 在其中是一个根系。由假设，存在 \(f \in V_{R}^{*}\) 使 \(\langle f, \alpha \rangle \neq 0\) 对一切 \(\alpha \in R\) 成立，且 \(\langle f, \mu \rangle = 0\) 。存在 \(R^{\vee}\) 的一个房 C 使 \(f \in C\) 。由第六章 §1 第 5 节 定理 2 (i)，存在 \(w \in W(R)\) 使 wf 属于对应于 B 的房，换言之，\(\langle wf, \alpha \rangle > 0\) 对一切 \(\alpha \in B\) 成立。记 \(w\mu = \sum_{\alpha \in B} t_{\alpha}\alpha\) 。则

\[
0 = \langle f, \mu \rangle = \langle w f, w \mu \rangle = \sum_ {\alpha \in B} t _ {\alpha} \langle w f, \alpha \rangle ,
\]

这证明某些 \(t_{\alpha}\) 为 \(> 0\) ，另一些为 \(< 0\) 。

引理 11. 设 \(\mu \in \mathrm{Q}\) 。若 \(\mu \notin \mathrm{R} \cup \{0\}\) ，则 \(\mathfrak{g}^{\mu} = 0\) 。若 \(\mu \in \mathrm{R}\) ，则 \(\dim \mathfrak{g}^{\mu} = 1\) 。 1) 若 \(\mu\) 不是 \(\mathrm{R}\) 中元素的倍数，则存在 \(w \in \mathrm{W}\) 使 \(w\mu \notin \mathrm{Q}_{+} \cup \mathrm{Q}_{-}\) （引理 10），于是 \(\mathfrak{a}^{w\mu} = 0\) ，\(\mathfrak{g}^{w\mu} = 0\) ，从而 \(\mathfrak{g}^{\mu} = 0\) （引理 9）。

\begin{enumerate}
\def\labelenumi{\arabic{enumi})}
\setcounter{enumi}{1}
\item
  设 \(\alpha \in \mathrm{B}\) ，\(m\) 是一个整数。由于 \(\mathfrak{a}_{+}\) 是以 \((x_{\alpha})_{\alpha \in \mathrm{B}}\) 为基础族的自由李代数，我们有 \(\dim \mathfrak{a}^{\alpha} = 1\) ，且 \(\mathfrak{a}^{m\alpha} = 0\) 对 \(m > 1\) 成立（第二章 §2 第 6 节 命题 4）。因此 \(\dim \mathfrak{g}^{\alpha} \leq 1\) ，且 \(\mathfrak{g}^{m\alpha} = 0\) 对 \(m > 1\) 成立。我们不可能有 \(\mathfrak{g}^{\alpha} = 0\) ，因为这将推出 \(x_{\alpha} \in \mathfrak{n} + \theta \mathfrak{n}\) ，从而 \(\mathfrak{n} + \theta \mathfrak{n}\) 包含 \(h_{\alpha} = -[x_{\alpha}, x_{-\alpha}]\) ，然而 \(\mathfrak{a}^0 \cap (\mathfrak{n} + \theta \mathfrak{n}) = 0\) 。于是，\(\dim \mathfrak{g}^{\alpha} = 1\) 。
\item
  若 \(\mu \in \mathrm{R}\) ，存在 \(w \in \mathrm{W}(\mathrm{R})\) 使 \(w(\mu) \in \mathrm{B}\) （第六章 §1 第 5 节 命题 15），故 \(\dim \mathfrak{g}^{\mu} = \dim \mathfrak{g}^{w\mu} = 1\) 。此外，若 \(n\) 是 \(>1\) 的整数，则 \(\mathfrak{g}^{nw(\mu)} = 0\) ，从而 \(\mathfrak{g}^{n\mu} = 0\) 。
\end{enumerate}

定理 1 的证明。

由于 \(\dim \mathfrak{g}^0 = \operatorname{CardB}\) ，由引理 11 推出 \(\mathfrak{g}\) 是有限维的，维数等于 \(\operatorname{CardB} + \operatorname{CardR}\) 。我们证明 \(\mathfrak{g}\) 是半单的。设 \(\mathfrak{k}\) 是 \(\mathfrak{g}\) 的一个交换理想。由于 \(\mathfrak{k}\) 在 \(\operatorname{ad}(\mathfrak{h})\) 之下稳定，\(\mathfrak{k} = (\mathfrak{k} \cap \mathfrak{h}) + \sum_{\mu \in \mathrm{R}} (\mathfrak{k} \cap \mathfrak{g}^\mu)\) 。

显然，对一切 \(\alpha \in \mathrm{B}\) ，\(\mathfrak{g}^{\alpha} + \mathfrak{g}^{-\alpha} + kH_{\alpha}\) 同构于 \(\mathfrak{sl}(2,k)\) 。鉴于引理 9，对一切 \(\mu \in \mathrm{R}\) ，\(\mathfrak{g}^{\mu}\) 包含于 \(\mathfrak{g}\) 的一个同构于 \(\mathfrak{sl}(2,k)\) 的子代数中；因此，\(\mathfrak{k} \cap \mathfrak{g}^{\mu} = 0\) ，故 \(\mathfrak{k} \subset \mathfrak{h}\) ；从而

\[
[ \mathfrak {k}, \mathfrak {g} ^ {\mu} ] \subset \mathfrak {k} \cap \mathfrak {g} ^ {\mu} = 0,
\]

故 \(\mu_{\mathrm{B}}(\mathfrak{k}) = 0\) 对一切 \(\mu \in \mathrm{R}\) 成立。由此推出 \(\mathfrak{k} = 0\) ，这就证明了 \(\mathfrak{g}\) 是半单的。

设 \(\mu \in \mathrm{R}\) 。存在 \(\alpha \in \mathrm{B}\) 使 \(\langle \mu, \alpha^{\vee} \rangle \neq 0\) ，于是 \((\operatorname{ad} H_{\alpha})|_{\mathfrak{g}^{\mu}}\) 是一个非零位似。因此，\(\mathfrak{h}\) 等于它在 \(\mathfrak{g}\) 中的正规化子，从而是 \(\mathfrak{g}\) 的一个嘉当子代数。对一切 \(u \in \mathfrak{h}\) ，\(\operatorname{ad} u\) 是可对角化的，故 \((\mathfrak{g}, \mathfrak{h})\) 是一个分裂半单李代数。

对一切 \(\mu \in \mathrm{R}\) ，显然 \(\mu_{\mathrm{B}}\) 是 \((\mathfrak{g},\mathfrak{h})\) 的根，且 \(\mathfrak{g}^{\mu}\) 是相应的特征子空间。\((\mathfrak{g},\mathfrak{h})\) 的根的数目是 \(\dim \mathfrak{g} - \dim \mathfrak{h} = \operatorname{Card}\mathrm{R}\) 。于是，映射 \(\mu \mapsto \mu_{\mathrm{B}}\) （从 V 到 \(\mathfrak{h}^*\) ）把 R 映为 \((\mathfrak{g},\mathfrak{h})\) 的根系。

\subsection{唯一性定理}

命题 4. 设 \((\mathfrak{g},\mathfrak{h},\mathrm{B},(X_{\alpha})_{\alpha\in\mathrm{B}})\) 是一个带标架的半单李代数。设 \((n(\alpha,\beta))_{\alpha,\beta\in\mathrm{B}}\) 与 \((X_{\alpha})_{\alpha\in\mathrm{B}\cup(-\mathrm{B})}\) 是相应的 Cartan 矩阵与生成族。

\begin{enumerate}
\def\labelenumi{(\roman{enumi})}
\item
  族 \(((X_{\alpha}, H_{\alpha}, X_{-\alpha}))_{\alpha \in \mathrm{B}}\) 连同第 3 节的关系式 (16) 至 (22) 构成 \(\mathfrak{g}\) 的一个呈现。
\item
  族 \((X_{\alpha})_{\alpha \in \mathrm{B}}\) 连同第 3 节的关系式 (21) 构成 \(\mathfrak{g}\) 的由 \((X_{\alpha})_{\alpha \in \mathrm{B}}\) 所生成的子代数的一个呈现。
\end{enumerate}

设 R 是 \((\mathfrak{g},\mathfrak{h})\) 的根系。把第 2 节与第 3 节中的构造应用于 R 与 B ，我们得到一些对象，将用 \(a', g', X'_{\alpha}, H'_{\alpha}, \ldots\) 来记它们，以代替 \(a, g, X_{\alpha}, H_{\alpha}, \ldots\)

存在同态 \(\varphi\) （从李代数 \(\mathfrak{g}'\) 到李代数 \(\mathfrak{g}\) ），使 \(\varphi(X_{\alpha}') = X_{\alpha}\) ，\(\varphi(H_{\alpha}') = H_{\alpha}\) ，\(\varphi(X_{-\alpha}') = X_{-\alpha}\) 对一切 \(\alpha \in B\) 成立（命题 1）。由于 \(\dim \mathfrak{g}' = \operatorname{Card} R + \operatorname{Card} B = \dim \mathfrak{g}\) ，\(\varphi\) 是双射。这就证明了 (i)。

\(\mathfrak{g}' = \mathfrak{a}' / (\mathfrak{n}'\oplus \theta'\mathfrak{n}') = (\mathfrak{a}_+^\prime \oplus \mathfrak{a}^{\prime 0}\oplus \mathfrak{a}_-^\prime) / (\mathfrak{n}'\oplus \theta'\mathfrak{n}')\) 中由 \((X_{\alpha}^{\prime})_{\alpha \in \mathrm{B}}\) 生成的子代数可与 \(\mathfrak{a}_+^\prime /\mathfrak{n}'\) 等同。鉴于命题 3 与 \(\mathfrak{n}'\) 的定义，这就证明了 (ii)。

推论. 每个带标架的半单李代数都可由一个带标架的半单 \(\mathbf{Q}\) -李代数经由基域从 \(\mathbf{Q}\) 到 \(k\) 的扩张而得到。

这由该命题立即得出。

定理 2. 设 \((\mathfrak{g},\mathfrak{h},\mathrm{B},(X_{\alpha})_{\alpha \in \mathrm{B}})\) 与 \((\mathfrak{g}',\mathfrak{h}',\mathrm{B}',(X'_{\alpha})_{\alpha \in \mathrm{B}'})\) 是带标架的半单李代数，设 R 与 R' 是 \((\mathfrak{g},\mathfrak{h})\) 与 \((\mathfrak{g}',\mathfrak{h}')\) 的根系，设 \((n(\alpha ,\beta))_{\alpha ,\beta \in \mathrm{B}}\) （相应地 \((n^{\prime}(\alpha ,\beta))_{\alpha ,\beta \in \mathrm{B}'}\) ）是 R （相应地 R' ）相对于 B （相应地 B' ）的 Cartan 矩阵，并设 \(\Delta\) （相应地 \(\Delta'\) ）是 R （相应地 R' ）相对于 B （相应地 B' ）的邓金图。

\begin{enumerate}
\def\labelenumi{(\roman{enumi})}
\item
  若 \(\varphi\) 是从 \(\mathfrak{h}^*\) 到 \(\mathfrak{h}'^*\) 的同构，使 \(\varphi(\mathrm{R}) = \mathrm{R}'\) 且 \(\varphi(\mathrm{B}) = \mathrm{B}'\) ，则存在唯一的同构 \(\psi\) （从 \((\mathfrak{g}, \mathfrak{h}, \mathrm{B}, (X_{\alpha})_{\alpha \in \mathrm{B}})\) 到 \((\mathfrak{g}', \mathfrak{h}', \mathrm{B}', (X_{\alpha}')_{\alpha \in \mathrm{B}'})\) ），使 \(\psi|_{\mathfrak{h}} = {}^t\varphi^{-1}\) 。
\item
  若 \(f\) 是从 \(\mathrm{B}\) 到 \(\mathrm{B}'\) 的双射，使 \(n'(f(\alpha), f(\beta)) = n(\alpha, \beta)\) 对一切 \(\alpha, \beta \in \mathrm{B}\) 成立，则存在从 \((\mathfrak{g}, \mathfrak{h}, \mathrm{B}, (X_{\alpha})_{\alpha \in \mathrm{B}})\) 到 \((\mathfrak{g}', \mathfrak{h}', \mathrm{B}', (X_{\alpha}')_{\alpha \in \mathrm{B}'}).\) 的同构
\item
  若存在从 \(\Delta\) 到 \(\Delta'\) 的同构，则存在从 \((\mathfrak{g},\mathfrak{h},\mathrm{B},(X_{\alpha})_{\alpha \in \mathrm{B}})\) 到 \((\mathfrak{g}',\mathfrak{h}',\mathrm{B}',(X_{\alpha}')_{\alpha \in \mathrm{B}'})\) 的同构。
\end{enumerate}

这由命题 4 (i) 立即得出（对部分 (iii) 要用到第六章 §4 第 2 节 命题 1）。

概要. 对每个可分裂半单李代数 g 都联系着一个邓金图，它在同构意义下确定 g （定理 2 (iii)）。此图非空且连通，当且仅当 g 是单的（§3 第 2 节 命题 6 的推论 1）。由第 3 节的定理 1 以及第六章 §4 第 2 节 定理 3，可分裂单李代数就是 \(A_{l}\) 型（ \(l \geq 1\) ）、\(B_{l}\) 型（ \(l \geq 2\) ）、\(C_{l}\) 型（ \(l \geq 3\) ）、\(D_{l}\) 型（ \(l \geq 4\) ）、\(E_{6}\) 型、\(E_{7}\) 型、\(E_{8}\) 型、\(F_{4}\) 型、\(G_{2}\) 型的李代数。此列表中任何两个代数都互不同构。

命题 5. 设 \((\mathfrak{g},\mathfrak{h},\mathrm{B},(X_{\alpha})_{\alpha \in \mathrm{B}})\) 是一个带标架的半单李代数，\((X_{\alpha})_{\alpha \in \mathrm{B}\cup (-\mathrm{B})}\) 是相应的生成族。存在唯一的自同构 \(\theta\) （\(\mathfrak{g}\) 的），使 \(\theta (X_{\alpha}) = X_{-\alpha}\) 对一切 \(\alpha \in \mathrm{B}\cup (-\mathrm{B})\) 成立。我们有 \(\theta^2 = \mathrm{Id}_{\mathfrak{g}}\) ，且 \(\theta (h) = -h\) 对一切 \(h\in \mathfrak{h}\) 成立。

唯一性是显然的，因为 \((X_{\alpha})_{\alpha\in\mathrm{B}\cup(-\mathrm{B})}\) 生成李代数 g。鉴于命题 4，\(\theta\) 的存在性由第 3 节中定理 1 之前所述的内容得出。

推论. 设 \((\mathfrak{g},\mathfrak{h})\) 是一个分裂半单李代数。则 \((\mathfrak{g},\mathfrak{h})\) 拥有一个谢瓦莱系（§2 第 4 节 定义 3）。

设 R 是 \((\mathfrak{g},\mathfrak{h})\) 的根系。对一切 \(\alpha\in R\) ，设 \(X_{\alpha}\) 是 \(g^{\alpha}\) 的一个非零元素。假设诸 \(X_{\alpha}\) 的选取使 \([X_{\alpha},X_{-\alpha}]=-H_{\alpha}\) 对一切 \(\alpha\in R\) 成立（§2 第 4 节 引理 2）。设 B 是 R 的一个基，\(\theta\) 是 g 的自同构，使 \(\theta(X_{\alpha})=X_{-\alpha}\) 对一切 \(\alpha\in\mathrm{B}\cup(-\mathrm{B})\) 成立。我们有 \(\theta|_{\mathfrak{h}}=-\mathrm{Id}_{\mathfrak{h}}\) 。于是，对一切 \(\alpha\in R\) ，存在 \(t_{\alpha}\in k^{*}\) 使 \(\theta X_{\alpha}=t_{\alpha}X_{-\alpha}\) 。我们有

\[
\begin{array}{c} t _ {\alpha} t _ {- \alpha} H _ {\alpha} = [ t _ {\alpha} X _ {- \alpha}, t _ {- \alpha} X _ {\alpha} ] = [ \theta X _ {\alpha}, \theta X _ {- \alpha} ] = \theta ([ X _ {\alpha}, X _ {- \alpha} ]) \\ = \theta (- H _ {\alpha}) = H _ {\alpha} \end{array}
\]

故 \(t_{\alpha}t_{-\alpha} = 1\) 对一切 \(\alpha \in \mathrm{R}\) 成立。按 §2 第 4 节引入诸 \(\mathrm{N}_{\alpha \beta}\) 。若 \(\alpha, \beta, \alpha + \beta \in \mathrm{R}\) ，

\[
\begin{array}{r} \mathrm{N} _ {- \alpha , - \beta} t _ {\alpha} t _ {\beta} X _ {- \alpha - \beta} = t _ {\alpha} t _ {\beta} [ X _ {- \alpha}, X _ {- \beta} ] = [ \theta X _ {\alpha}, \theta X _ {\beta} ] = \theta ([ X _ {\alpha}, X _ {\beta} ]) \\ = \mathrm{N} _ {\alpha \beta} \theta X _ {\alpha + \beta} = \mathrm{N} _ {\alpha \beta} t _ {\alpha + \beta} X _ {- \alpha - \beta} \end{array}
\]

于是，鉴于 §2 第 4 节 引理 4，

\[
(q + 1) ^ {2} t _ {\alpha} t _ {\beta} = \mathrm{N} _ {\alpha \beta} ^ {2} t _ {\alpha + \beta}
\]

其中 \(q\) 是一个整数。由此推出，若 \(t_{\alpha}\) 与 \(t_{\beta}\) 是 \(k^*\) 中的平方元，则 \(t_{\alpha + \beta}\) 也是。由于 \(t_{\alpha} = 1\) 对一切 \(\alpha \in \mathrm{B}\) 成立，第六章 §1 第 6 节 命题 19 证明 \(t_{\alpha}\) 对一切 \(\alpha \in \mathrm{R}\) 都是平方元。对一切 \(\alpha \in \mathrm{R}\) ，选取 \(u_{\alpha} \in k\) 使 \(u_{\alpha}^{2} = t_{\alpha}\) 。这一选取可以做到使 \(u_{\alpha}u_{-\alpha} = 1\) 对一切 \(\alpha \in \mathrm{R}\) 成立。置 \(X_{\alpha}' = u_{\alpha}^{-1}X_{\alpha}\) 。则对一切 \(\alpha \in \mathrm{R}\) ，

\[
X _ {\alpha} ^ {\prime} \in \mathfrak {g} ^ {\alpha}, [ X _ {\alpha} ^ {\prime}, X _ {- \alpha} ^ {\prime} ] = [ X _ {\alpha}, X _ {- \alpha} ] = - H _ {\alpha},
\]

且 \(\theta(X_{\alpha}^{\prime}) = \theta(u_{\alpha}^{-1}X_{\alpha}) = u_{\alpha}^{-1}t_{\alpha}X_{-\alpha} = u_{\alpha}X_{-\alpha} = u_{\alpha}u_{-\alpha}X_{-\alpha}^{\prime} = X_{-\alpha}^{\prime}\) ，于是 \((X_{\alpha}^{\prime})_{\alpha \in \mathrm{R}}\) 是 \((\mathfrak{g},\mathfrak{h})\) 的一个谢瓦莱系。

\section{半单李代数的自同构}

在本节中，g 表示一个半单李代数。

\subsection{带标架的半单李代数的自同构}

回顾（第七章 §3 第 1 节），\(\operatorname{Aut}(\mathfrak{g})\) 表示 \(\mathfrak{g}\) 的自同构群。若 \(\mathfrak{h}\) 是 \(\mathfrak{g}\) 的一个嘉当子代数，我们用 \(\operatorname{Aut}(\mathfrak{g},\mathfrak{h})\) 表示 \(\mathfrak{g}\) 的使 \(\mathfrak{h}\) 稳定的自同构组成的群。假设 \(\mathfrak{h}\) 是分裂的，并设 R 是 \((\mathfrak{g},\mathfrak{h})\) 的根系。若 \(s\in\operatorname{Aut}(\mathfrak{g},\mathfrak{h})\) ，则 \(s|_{\mathfrak{h}}\) 的反步映射是 A(R) （R 的自同构群）的一个元素，我们在本节中将它记作 \(\varepsilon(s)\) 。于是

\[
\varepsilon : \operatorname{Aut} (\mathfrak {g}, \mathfrak {h}) \to \mathrm{A} (\mathrm{R})
\]

是一个群同态。

对任何根系 R 及 R 的任何基 B ，我们用 \(\operatorname{Aut}(R,B)\) 表示使 B 稳定的 R 的自同构组成的群。回顾（第六章 §1 第 5 节 命题 16 与 §4 第 2 节 命题 1 的推论），A(R) 是 \(\operatorname{Aut}(R,B)\) 与 W(R) 的半直积，且 A(R)/W(R) 典则同构于 R 的邓金图的自同构群。

命题 1. 设 \((\mathfrak{g},\mathfrak{h},\mathrm{B},(X_{\alpha})_{\alpha \in \mathrm{B}})\) 是一个带标架的半单李代数，R 是 \((\mathfrak{g},\mathfrak{h})\) 的根系。设 G 是这样的 \(s\in \operatorname {Aut}(\mathfrak{g},\mathfrak{h})\) 组成的集合：它们使 B 稳定，且使 \(s(X_{\alpha}) = X_{\varepsilon (s)\alpha}\) 对一切 \(\alpha \in \mathrm{B}\) 成立（换言之，即 \((\mathfrak{g},\mathfrak{h},\mathrm{B},(X_{\alpha})_{\alpha \in \mathrm{B}})\) 的自同构组成的集合）。则 \(\varepsilon\) 在 \(\mathrm{G}\) 上的限制是从 \(\mathrm{G}\) 到 \(\operatorname {Aut}(\mathrm{R},\mathrm{B})\) 的同构。

若 \(s \in G\) ，显然 \(\varepsilon(s) \in \operatorname{Aut}(\mathrm{R}, \mathrm{B})\) 。另一方面，映射

\[
\varepsilon|_{\mathrm{G}}: \mathrm{G} \rightarrow \operatorname{Aut} (\mathrm{R}, \mathrm{B})
\]

由 §4 第 4 节 定理 2 是双射的。

\subsection{分裂半单李代数的自同构}

设 E 是一个交换群，\(A = \bigoplus_{\gamma \in E} A^{\gamma}\) 是一个 E-分次代数。对任何同态 \(\varphi\) （从群 E 到乘法群 \(k^{*}\) 的），用 \(f(\varphi)\) 表示从 A 到 A 的 k-线性映射，它在每个 \(A^{\gamma}\) 上的限制是比为 \(\varphi(\gamma)\) 的位似；显然 \(f(\varphi)\) 是分次代数 A 的自同构，且 f 是从群 \(\operatorname{Hom}(E, k^{*})\) 到分次代数 A 的自同构群的同态。

设 h 是 g 的一个分裂嘉当子代数，R 是 \((\mathfrak{g}, \mathfrak{h})\) 的根系。回顾 P(R) （相应地 Q(R) ）表示 R 的权格（相应地根权格）。置

\[
\mathrm{T} _ {\mathrm{P}} = \operatorname{Hom} (\mathrm{P} (\mathrm{R}), k ^ {*}) \quad \mathrm{T} _ {\mathrm{Q}} = \operatorname{Hom} (\mathrm{Q} (\mathrm{R}), k ^ {*}).
\]

我们可以把 \(g = g^{0} + \sum_{\alpha \in R} g^{\alpha}\) 视为一个 \(Q(R)\) -分次代数。前面的注记定义了一个从 \(T_{Q}\) 到 \(\operatorname{Aut}(\mathfrak{g}, \mathfrak{h})\) 的典则同态，它在本节中将记作 f。另一方面，从 \(Q(R)\) 到 \(P(R)\) 的典则单射定义了一个从 \(T_{P}\) 到 \(T_{Q}\) 的同态，将记作 q：

\[
\mathrm{T} _ {\mathrm{P}} \quad \stackrel {{q}} {{\longrightarrow}} \quad \mathrm{T} _ {\mathrm{Q}} \quad \stackrel {{f}} {{\longrightarrow}} \quad \operatorname{Aut} (\mathfrak {g}, \mathfrak {h}).
\]

若 \(s \in \operatorname{Aut}(\mathfrak{g}, \mathfrak{h})\) ，用 \(s^*\) 表示 \({}^t(s|_{\mathfrak{h}})^{-1}\) 在 \(\mathrm{Q}(\mathrm{R})\) 上的限制。则对一切 \(\varphi \in \mathrm{T}_{\mathrm{Q}}\) ，

\[
f (\varphi \circ s ^ {*}) = s ^ {- 1} \circ f (\varphi) \circ s.\tag{1}
\]

事实上，设 \(\gamma \in \mathrm{Q}(\mathrm{R})\) 与 \(x\in \mathfrak{g}^{\gamma}\) ；则 \(sx\in \mathfrak{g}^{s^{*}\gamma}\) ，且

\[
f (\varphi \circ s ^ {*}) x = (\varphi \circ s ^ {*}) (\gamma). x = s ^ {- 1} (\varphi (s ^ {*} \gamma) s x) = (s ^ {- 1} \circ f (\varphi) \circ s) (x).
\]

命题 2. 同态序列

\[
1 \longrightarrow \mathrm {T_ {Q}} \stackrel {f} {\longrightarrow} \mathrm{Aut} (\mathfrak {g}, \mathfrak {h}) \stackrel {\varepsilon} {\longrightarrow} \mathrm{A(R)} \longrightarrow 1
\]

是正合的。

\begin{enumerate}
\def\labelenumi{\alph{enumi})}
\item
  设 \(\varphi \in \operatorname{Ker} f\) 。则 \(\varphi(\alpha) = 1\) 对一切 \(\alpha \in \mathrm{R}\) 成立。由于 \(\mathrm{R}\) 生成群 \(\mathrm{Q}(\mathrm{R})\) ，\(\varphi\) 是 \(\mathrm{T}_{\mathrm{Q}}\) 的单位元。
\item
  设 \(\varphi \in \mathrm{T}_{\mathrm{Q}}\) 。\(f(\varphi)\) 在 \(\mathfrak{h} = \mathfrak{g}^0\) 上的限制是恒等映射，故
\end{enumerate}

\(\operatorname{Im} f \subset \operatorname{Ker} \varepsilon\) 。

\begin{enumerate}
\def\labelenumi{\alph{enumi})}
\setcounter{enumi}{2}
\tightlist
\item
  设 \(s \in \operatorname{Ker} \varepsilon\) 。则 \(s|_{\mathfrak{h}} = \mathrm{Id}_{\mathfrak{h}}\) 。对一切 \(\alpha \in R\) ，我们有 \(s(\mathfrak{g}^{\alpha}) = \mathfrak{g}^{\alpha}\) ，且存在 \(t_{\alpha} \in k^{*}\) 使 \(sx = t_{\alpha}x\) 对一切 \(x \in g^{\alpha}\) 成立。写下 \(s \in \operatorname{Aut}(\mathfrak{g})\) 这一条件，我们得到诸关系式
\end{enumerate}

\[
t _ {\alpha} t _ {- \alpha} = 1 \qquad \text{对一切 } \alpha \in \mathrm{R}
\]

\[
t _ {\alpha} t _ {\beta} = t _ {\alpha + \beta} \quad \text { when } \alpha , \beta , \alpha + \beta \in \mathrm{R}.
\]

在这些条件下，存在 \(\varphi \in \mathrm{T}_{\mathrm{Q}}\) 使 \(\varphi (\alpha) = t_{\alpha}\) 对一切 \(\alpha \in \mathrm{R}\) 成立（第六章 §1 第 6 节 命题 19 的推论 2）。则 \(s = f(\varphi)\) 。因此，\(\operatorname{Ker} \varepsilon \subset \operatorname{Im} f\) 。

\begin{enumerate}
\def\labelenumi{\alph{enumi})}
\setcounter{enumi}{3}
\tightlist
\item
  \(\operatorname{Aut}(\mathfrak{g},\mathfrak{h})\) 在 \(\varepsilon\) 下的像由 §2 第 2 节 定理 2 的推论知包含 \(\mathrm{W}(\mathrm{R})\) ，并由命题 1 知包含 \(\operatorname{Aut}(\mathrm{R},\mathrm{B})\) 。因此这个像等于 \(\mathrm{A}(\mathrm{R})\) 。
\end{enumerate}

推论 1. 设 \((\mathrm{B}, (X_{\alpha})_{\alpha \in \mathrm{B}})\) 是 \((\mathfrak{g}, \mathfrak{h})\) 的一个标架。设 \(\mathrm{G}\) 是使该标架不变的 \(s \in \operatorname{Aut}(\mathfrak{g}, \mathfrak{h})\) 组成的集合。则 \(\operatorname{Aut}(\mathfrak{g}, \mathfrak{h})\) 是 \(\mathrm{G}\) 与 \(\varepsilon^{-1}(\mathrm{W}(\mathrm{R}))\) 的半直积。

事实上，由命题 1，\(\mathrm{G} \cap \varepsilon^{-1}(\mathrm{W}(\mathrm{R})) = \{1\}\) ，且

\[
\operatorname{Aut} (\mathfrak {g}, \mathfrak {h}) = \mathrm{G}. \varepsilon^ {- 1} (\mathrm{W(R)})
\]

因为 \(\varepsilon\) 是满射（命题 2）。

推论 2. 群 \(\varepsilon^{-1}(\mathrm{W}(\mathrm{R}))\) 单可递地作用在 \((\mathfrak{g},\mathfrak{h})\) 的标架组成的集合上。

事实上，由 §4 第 4 节 定理 2，\(\operatorname{Aut}(\mathfrak{g},\mathfrak{h})\) 可递地作用在 \((\mathfrak{g},\mathfrak{h})\) 的标架组成的集合上。推论 2 于是由推论 1 得出。

推论 3. 设 B 是 R 的一个基。群 \(\operatorname{Ker} \varepsilon = f(\mathrm{T}_{\mathrm{Q}})\) 单可递地作用在 \((\mathfrak{g}, \mathfrak{h})\) 的形如 \((\mathrm{B}, (X_{\alpha})_{\alpha \in \mathrm{B}})\) 的标架组成的集合上。

这立即由命题 2 得出。

设 \(\alpha \in \mathrm{R}\) ，\(X_{\alpha} \in \mathfrak{g}^{\alpha}\) ，\(X_{-\alpha} \in \mathfrak{g}^{-\alpha}\) 满足 \([X_{\alpha}, X_{-\alpha}] = -H_{\alpha}\) 。我们已看到（§2 第 2 节 定理 2），对一切 \(t \in k^{*}\) ，初等自同构

\[
\theta_ {\alpha} (t) = e ^ {\mathrm{ad} t X _ {\alpha}} e ^ {\mathrm{ad} t ^ {- 1} X _ {- \alpha}} e ^ {\mathrm{ad} t X _ {\alpha}}
\]

在 \(\mathfrak{h}\) 上的限制是 \(s_{\alpha}\) 的转置；于是 \(\varepsilon(\theta_{\alpha}(t)) = s_{\alpha}\) ，从而 \(\theta_{\alpha}(t)\theta_{\alpha}(-1) \in \operatorname{Ker} \varepsilon\) 。

引理 1. 设 \(\alpha \in \mathrm{R}\) 与 \(t\in k^{*}\) 。设 \(\varphi\) 是同态 \(\lambda \mapsto t^{\lambda (H_{\alpha})}\) （从 \(\mathrm{Q(R)}\) 到 \(k^*\) 的）。则 \(f(\varphi) = \theta_{\alpha}(t)\theta_{\alpha}(-1)\) 。

设 \(\rho\) 是 \(\mathfrak{sl}(2,k)\) 在 \(\mathfrak{g}\) 上的与 \(X_{\alpha}\) 相伴的表示。设 \(\pi\) 是 \(\mathbf{SL}(2,k)\) 的与 \(\rho\) 相容的表示。采用 §1 第 5 节 的记号 \(\theta(t), h(t)\) 。由于 \(\rho(H) = \operatorname{ad} H_{\alpha}\) ，\(\mathfrak{g}^{\lambda}\) 的元素的权是 \(\lambda(H_{\alpha})\) （对 \(\rho\) 而言）。由 §2 第 2 节，\(\theta_{\alpha}(t)\theta_{\alpha}(-1) = \pi(\theta(t)\theta(-1)) = \pi(h(t))\) 。因此 \(\theta_{\alpha}(t)\theta_{\alpha}(-1)\) 在 \(\mathfrak{g}^{\lambda}\) 上的限制是比为 \(t^{\lambda(H_{\alpha})}\) 的位似（§1 第 5 节 命题 6），由此得引理。

命题 3. 复合同态

\[
\mathrm{T} _ {\mathrm{P}} \quad \xrightarrow {q} \quad \mathrm{T} _ {\mathrm{Q}} \quad \xrightarrow {f} \quad \operatorname{Aut} (\mathfrak {g}, \mathfrak {h})
\]

的像包含于 \(\mathrm{Aut}_e(\mathfrak{g})\) 。

设 B 是 R 的一个基。则 \((H_{\alpha})_{\alpha\in\mathrm{B}}\) 是 \(R^{\vee}\) 的一个基，且 \((H_{\alpha})_{\alpha\in\mathrm{B}}\) 在 \(h^{*}\) 中的对偶基是群 P(R) 的一个基。因此群 \(T_{P}\) 由诸同态 \(\lambda\mapsto t^{\lambda(H_{\alpha})}\;(t\in k^{*},\alpha\in\mathrm{B})\) 生成。若 \(\varphi\) 是这样一个同态到 Q(R) 上的限制，引理 1 证明 \(f(\varphi)\in\operatorname{Aut}_{e}(\mathfrak{g})\) ，由此得命题。

设 \(\bar{k}\) 是 \(k\) 的一个代数闭包。把每个自同构 \(s\) （\(\mathfrak{g}\) 的）对应到自同构 \(s \otimes 1\) （\(\mathfrak{g} \otimes_k \bar{k}\) 的）的映射是从 \(\mathrm{Aut}(\mathfrak{g})\) 到 \(\mathrm{Aut}(\mathfrak{g} \otimes_k \bar{k})\) 的单同态。我们用 \(\mathrm{Aut}_0(\mathfrak{g})\) 表示 \(\mathrm{Aut}(\mathfrak{g})\) 的正规子群，它是 \(\mathrm{Aut}_e(\mathfrak{g} \otimes_k \bar{k})\) 在这个同态下的逆像；这就是 \(\mathfrak{g}\) 的那样一些自同构组成的集合：它们在把基域从 \(k\) 扩张到 \(\bar{k}\) 后成为初等的。显然 \(\mathrm{Aut}_e(\mathfrak{g})\) 与 \(\bar{k}\) 的选取无关，且 \(\mathrm{Aut}_e(\mathfrak{g}) \subset \mathrm{Aut}_0(\mathfrak{g})\) 。群 \(\mathrm{Aut}_0(\mathfrak{g})\) 与 \(\mathrm{Aut}_e(\mathfrak{g})\) 可以不同（第七章 §13 第 1 节）。若 \(\mathfrak{h}\) 是 \(\mathfrak{g}\) 的一个嘉当子代数，置

\[
\operatorname{Aut} _ {e} (\mathfrak {g}, \mathfrak {h}) = \operatorname{Aut} _ {e} (\mathfrak {g}) \cap \operatorname{Aut} (\mathfrak {g}, \mathfrak {h}), \quad \operatorname{Aut} _ {0} (\mathfrak {g}, \mathfrak {h}) = \operatorname{Aut} _ {0} (\mathfrak {g}) \cap \operatorname{Aut} (\mathfrak {g}, \mathfrak {h}).
\]

引理 2. 设 \(\mathfrak{h}\) 是 \(\mathfrak{g}\) 的一个分裂嘉当子代数，且 \(s \in \operatorname{Aut}_0(\mathfrak{g}, \mathfrak{h})\) 。假设 \(s\) 在 \(\sum_{\alpha \in \mathrm{R}} \mathfrak{g}^\alpha\) 上的限制不以 1 为特征值。则 \(\varepsilon(s) = 1\) 。

通过扩张 \(k\) ，我们化归到 \(s \in \mathrm{Aut}_e(\mathfrak{g}, \mathfrak{h})\) 的情形。\(s - 1\) 的幂零空间的维数至少是 \(\dim \mathfrak{h}\) （第七章 §4 第 4 节 命题 9）。因此 \((s - 1)|_{\mathfrak{h}}\) 是幂零的。由于 \(s|_{\mathfrak{h}} \in \mathrm{A}(\mathrm{R}^\vee)\) ，\(s|_{\mathfrak{h}}\) 是有限阶的，从而是半单的（第五章附录命题 2）。于是 \((s - 1)|_{\mathfrak{h}} = 0\) ，这就证明了 \(\varepsilon(s) = 1\) 。

引理 3. (i) 设 \(m = (\mathrm{P}(\mathrm{R}) : \mathrm{Q}(\mathrm{R}))\) 。若 \(\varphi\) 是某个元素的 \(m\) 次幂（\(\mathrm{T}_{\mathrm{Q}}\) 中的元素），则 \(\varphi \in q(\mathrm{T}_{\mathrm{P}})\) 。

\begin{enumerate}
\def\labelenumi{(\roman{enumi})}
\setcounter{enumi}{1}
\tightlist
\item
  若 \(k\) 是代数闭的，则 \(q(\mathrm{T}_{\mathrm{P}}) = \mathrm{T}_{\mathrm{Q}}\) 。
\end{enumerate}

存在 P(R) 的一个基 \((\lambda_{1},\ldots,\lambda_{l})\) 以及整数 \(n_{1}\geq1,\ldots,n_{l}\geq1\) ，使 \((n_{1}\lambda_{1},\ldots,n_{l}\lambda_{l})\) 是 Q(R) 的一个基。我们有 \(m=n_{1}\ldots n_{l}\) 。设 \(\psi\in T_{Q}\) ，并置 \(\psi(n_{1}\lambda_{1})=t_{1},\ldots,\psi(n_{l}\lambda_{l})=t_{l}\) 。对 \(i=1,\ldots,l\) ，置 \(m_{i}=\) \(\prod_{j \neq i} n_j\) 。设 \(\chi\) 是 \(T_P\) 的元素，使 \(\chi(\lambda_1) = t_1^{m_1}, \ldots, \chi(\lambda_l) = t_l^{m_l}\) 。则

\[
\chi (n _ {i} \lambda_ {i}) = t _ {i} ^ {m _ {i} n _ {i}} = t _ {i} ^ {m} = (\psi^ {m}) (n _ {i} \lambda_ {i})
\]

于是 \(\chi|_{\mathrm{Q}(\mathrm{R})}=\psi^{m}\) 。这就证明了 (i)。若 k 是代数闭的，\(k^{*}\) 的每个元素都是 \(k^{*}\) 的某个元素的 m 次幂，故 \(T_{Q}\) 的每个元素都是 \(T_{Q}\) 的某个元素的 m 次幂；于是 (ii) 由 (i) 得出。

命题 4. 我们有 \(f(\mathrm{T}_{\mathrm{Q}}) \subset \mathrm{Aut}_0(\mathfrak{g}, \mathfrak{h})\) 与 \(\varepsilon^{-1}(\mathrm{W}(\mathrm{R})) = \mathrm{Aut}_0(\mathfrak{g}, \mathfrak{h})\) 。

\begin{enumerate}
\def\labelenumi{\alph{enumi})}
\tightlist
\item
  设 \(\varphi \in \mathrm{T}_{\mathrm{Q}}\) ，并设 \(\bar{k}\) 是 \(k\) 的一个代数闭包。由引理 3，\(\varphi\) 可以扩张为 \(\operatorname{Hom}(\mathrm{P}(\mathrm{R}), k^*)\) 的一个元素。由命题 3，
\end{enumerate}

\[
f (\varphi) \otimes 1 \in \operatorname{Aut} _ {e} (\mathfrak {g} \otimes_ {k} \bar {k}, \mathfrak {h} \otimes_ {k} \bar {k}).
\]

因此 \(f(\varphi) \in \mathrm{Aut}_0(\mathfrak{g},\mathfrak{h})\) ，且 \(\operatorname{Ker} \varepsilon \subset \mathrm{Aut}_0(\mathfrak{g},\mathfrak{h})\) 。

\begin{enumerate}
\def\labelenumi{\alph{enumi})}
\setcounter{enumi}{1}
\item
  \(\mathrm{Aut}_e(\mathfrak{g},\mathfrak{h})\) 在 \(\varepsilon\) 下的像包含 \(\mathrm{W(R)}\) （§2 第 2 节 定理 2 的推论）。鉴于 a)，我们看到 \(\varepsilon^{-1}(\mathrm{W(R)})\subset \mathrm{Aut}_0(\mathfrak{g},\mathfrak{h})\) 。
\item
  剩下要证 \(\mathrm{Aut}_0(\mathfrak{g},\mathfrak{h})\subset \varepsilon^{-1}(\mathrm{W}(\mathrm{R}))\) 。鉴于 b)，只需证明 \(\varepsilon (\mathrm{Aut}_0(\mathfrak{g},\mathfrak{h}))\cap \mathrm{Aut}(\mathrm{R},\mathrm{B})\) （其中 B 表示 R 的一个基）约化为 \{1\}。
\end{enumerate}

设 \(s \in \mathrm{Aut}_0(\mathfrak{g}, \mathfrak{h})\) 满足 \(\varepsilon(s) \in \mathrm{Aut}(\mathrm{R}, \mathrm{B})\) 。\(\mathrm{A}(\mathrm{R})\) 的由 \(\varepsilon(s)\) 生成的子群在 \(\mathrm{R}\) 上只有有限个轨道。设 \(\mathrm{U}\) 是这样一个轨道，其基数是 \(r\) ，并设 \(\mathfrak{g}^{\mathrm{U}} = \sum_{\beta \in \mathrm{U}} \mathfrak{g}^{\beta}\) 。设 \(\beta_1 \in \mathrm{U}\) ，并置 \(\beta_i = \varepsilon(s)^{i-1}\beta_1\) 对 \(1 \leq i \leq r\) ，于是 \(\mathrm{U} = \{\beta_1, \ldots, \beta_r\}\) 。设 \(X_{\beta_1}\) 是 \(\mathfrak{g}^{\beta_1}\) 的一个非零元素，并置 \(X_{\beta_i} = s^{i-1}X_{\beta_1}\) 对 \(1 \leq i \leq r\) 。存在 \(c_{\mathrm{U}} \in k^*\) 使 \(s^r X_{\beta_1} = c_{\mathrm{U}}X_{\beta_1}\) ，于是 \(s^r X_{\beta_i} = c_{\mathrm{U}}X_{\beta_i}\) 对一切 \(i\) 成立，从而 \(s^r |_{\mathfrak{g}^{\mathrm{U}}} = c_{\mathrm{U}}.1\) 。设 \(\varphi \in T_{\mathrm{Q}}\) ，并设 \(s' = s \circ f(\varphi)\) ，由 \(a)\) 它是 \(\mathrm{Aut}_0(\mathfrak{g}, \mathfrak{h})\) 的元素。我们有 \(s'^r |_{\mathfrak{g}^{\mathrm{U}}} = c_{\mathrm{U}}'.1\) ，其中

\[
c _ {\mathrm{U}} ^ {\prime} = c _ {\mathrm{U}} \prod_ {i = 1} ^ {r} \varphi (\beta_ {i}) = c _ {\mathrm{U}} \varphi \left(\sum_ {i = 1} ^ {r} \beta_ {i}\right).
\]

置 \(\mathrm{B} = \{\alpha_1, \ldots, \alpha_l\}\) 及 \(\sum_{i=1}^{r} \beta_i = \sum_{j=1}^{l} m_j^{\mathrm{U}} \alpha_j\) 。由于 \(\varepsilon(s) \in \mathrm{Aut}(\mathrm{R}, \mathrm{B})\) ，诸 \(m_j^{\mathrm{U}}\) 是同号且不全为零的整数。我们有

\[
c _ {\mathrm{U}} ^ {\prime} = c _ {\mathrm{U}} \prod_ {j = 1} ^ {l} \varphi (\alpha_ {j}) ^ {m _ {j} ^ {\mathrm{U}}}.
\]

现在可以选取 \(\varphi\) ，使对每个轨道 U 都有 \(c_{\mathrm{U}}^{\prime} \neq 1\) ；事实上，这归结为选取元素 \(\varphi(\alpha_1) = t_1, \ldots, \varphi(\alpha_l) = t_l\) （\(k^*\) 中的），它们不被有限个 \(t_1, \ldots, t_l\) 的不恒为零的多项式所零化。对这样的 \(\varphi\) 的选取，由引理 2 有 \(\varepsilon(s') = 1\) ，故

\[
\varepsilon (s) = \varepsilon (s ^ {\prime}) \varepsilon (f (\varphi)) ^ {- 1} = 1.
\]

推论. 设 B 是 R 的一个基。群 \(\operatorname{Aut}(\mathfrak{g},\mathfrak{h})\) 同构于群 \(\operatorname{Aut}(\mathrm{R},\mathrm{B})\) 与 \(\operatorname{Aut}_{0}(\mathfrak{g},\mathfrak{h})\) 的半直积。

这由命题 1、命题 2 的推论 1 及命题 4 得出。

注. 设 \(\varepsilon^{\prime},\varepsilon^{\prime\prime}\) 是 \(\varepsilon\) 到 \(\operatorname{Aut}_{0}(\mathfrak{g},\mathfrak{h}),\operatorname{Aut}_{e}(\mathfrak{g},\mathfrak{h})\) 上的限制。设 \(f^{\prime}\) 是从 \(T_{P}\) 到 \(\operatorname{Aut}_{e}(\mathfrak{g},\mathfrak{h})\) 的、由 f 经由从 \(Q(R)\) 到 \(P(R)\) 的典则单射诱导的同态。在前面我们已建立了如下的交换图：

\[
\begin{array}{c c c c c c c c c} 1 & \longrightarrow & \mathrm{T} _ {\mathrm{Q}} & \stackrel {{f}} {{\longrightarrow}} & \mathrm{Aut} (\mathfrak {g}, \mathfrak {h}) & \stackrel {{\varepsilon}} {{\longrightarrow}} & \mathrm{A} (\mathrm{R}) & \longrightarrow & 1 \\ & & \uparrow & & \uparrow & & \uparrow & & \\ 1 & \longrightarrow & \mathrm{T} _ {\mathrm{Q}} & \stackrel {{f}} {{\longrightarrow}} & \mathrm{Aut} _ {0} (\mathfrak {g}, \mathfrak {h}) & \stackrel {{\varepsilon^ {\prime}}} {{\longrightarrow}} & \mathrm{W} (\mathrm{R}) & \longrightarrow & 1 \\ & & \uparrow_ {q} & & \uparrow & & \uparrow & & \\ & & \mathrm{T} _ {\mathrm{P}} & \stackrel {{f ^ {\prime}}} {{\longrightarrow}} & \mathrm{Aut} _ {e} (\mathfrak {g}, \mathfrak {h}) & \stackrel {{\varepsilon^ {\prime \prime}}} {{\longrightarrow}} & \mathrm{W} (\mathrm{R}) & \longrightarrow & 1 \end{array}
\]

其中除 q 外的竖直箭头表示典则单射。我们已看到（命题 2 与命题 4），前两行是正合的。在第三行中，同态 \(\varepsilon''\) 是满射（§2 第 2 节 定理 2 的推论）；可以证明其核是 \(f'(\mathrm{T}_{\mathrm{P}})\) （§7 习题 26 d)）。

\subsection{可分裂半单李代数的自同构}

命题 5. 假设 g 是可分裂的。群 \(\operatorname{Aut}_{0}(\mathfrak{g})\) 单可递地作用在 g 的标架组成的集合上。

设 \(e_1 = (\mathfrak{g}, \mathfrak{h}_1, \mathrm{B}_1, (X_\alpha^1)_{\alpha \in \mathrm{B}_1})\) ，\(e_2 = (\mathfrak{g}, \mathfrak{h}_2, \mathrm{B}_2, (X_\alpha^2)_{\alpha \in \mathrm{B}_2})\) 是 \(\mathfrak{g}\) 的两个标架。存在 \(\operatorname{Aut}_0(\mathfrak{g})\) 的至少一个元素把 \(e_1\) 变为 \(e_2\) （命题 1 与命题 4）。设 \(\bar{k}\) 是 \(k\) 的一个代数闭包。存在 \(\operatorname{Aut}_e(\mathfrak{g} \otimes_k \bar{k})\) 的一个元素把 \(\mathfrak{h}_1 \otimes_k \bar{k}\) 变为 \(\mathfrak{h}_2 \otimes_k \bar{k}\) （第七章 §3 第 2 节 定理 1）。于是，由命题 4 与命题 2 的推论 2，存在元素 \(\varphi\) （\(\operatorname{Aut}_e(\mathfrak{g} \otimes_k \bar{k})\) 的），它把标架 \((\mathfrak{g} \otimes_k \bar{k}, \mathfrak{h}_1 \otimes_k \bar{k}, \mathrm{B}_1, (X_\alpha^1)_{\alpha \in \mathrm{B}_1})\) （\(\mathfrak{g} \otimes_k \bar{k}\) 的）变为标架 \((\mathfrak{g} \otimes_k \bar{k}, \mathfrak{h}_2 \otimes_k \bar{k}, \mathrm{B}_2, (X_\alpha^2)_{\alpha \in \mathrm{B}_2})\) 。由于 \(\mathfrak{h}_1\) 与诸 \(X_\alpha^1\) （相应地 \(\mathfrak{h}_2\) 与诸 \(X_\alpha^2\) ）生成 \(\mathfrak{g}_1\) （相应地 \(\mathfrak{g}_2\) ），我们有 \(\varphi(\mathfrak{g}_1) = \mathfrak{g}_2\) ，故 \(\varphi\) 形如 \(\psi \otimes 1\) ，其中 \(\psi \in \operatorname{Aut}_0(\mathfrak{g})\) ，且 \(\psi\) 把 \(e_1\) 变为 \(e_2\) 。

推论 1. 设 \((\mathfrak{g},\mathfrak{h},\mathrm{B},(X_{\alpha})_{\alpha\in\mathrm{B}})\) 是 g 的一个标架，G 是使该标架不变的 g 的自同构组成的群（同构于 \(\operatorname{Aut}(\mathrm{R},\mathrm{B})\) ）。则 \(\operatorname{Aut}(\mathfrak{g})\) 是 G 与 \(\operatorname{Aut}_{0}(\mathfrak{g})\) 的半直积。

事实上，\(\operatorname{Aut}(\mathfrak{g})\) 的每个元素把 \((\mathfrak{g},\mathfrak{h},\mathrm{B},(X_{\alpha})_{\alpha\in\mathrm{B}})\) 变为 g 的一个标架。由命题 5，\(\operatorname{Aut}(\mathfrak{g})\) 模 \(\operatorname{Aut}_{0}(\mathfrak{g})\) 的每个陪集与 G 恰好交于一点。证毕。

由推论 1 可知，群 \(\operatorname{Aut}(\mathfrak{g})/\operatorname{Aut}_{0}(\mathfrak{g})\) 可以与 \(\operatorname{Aut}(\mathrm{R},\mathrm{B})\) 等同，并且同构于 R 的邓金图的自同构群。

推论 2. \(\operatorname{Aut}(\mathfrak{g}) = \operatorname{Aut}_0(\mathfrak{g})\) ，当 \(\mathfrak{g}\) 是 \(\mathrm{A}_1\) , \(\mathrm{B}_n\) ( \(n \geq 2\) ), \(\mathrm{C}_n\) ( \(n \geq 2\) ), \(\mathrm{E}_7\) , \(\mathrm{E}_8\) , \(\mathrm{F}_4\) , \(\mathrm{G}_2\) 型的可分裂单李代数时。商群 \(\operatorname{Aut}(\mathfrak{g}) / \operatorname{Aut}_0(\mathfrak{g})\) 的阶是 2 ，当 \(\mathfrak{g}\) 是 \(\mathrm{A}_n\) ( \(n \geq 2\) ), \(\mathrm{D}_n\) ( \(n \geq 5\) ), \(\mathrm{E}_6\) 型时；它同构于 \(\mathfrak{S}_3\) ，当 \(\mathfrak{g}\) 是 \(\mathrm{D}_4\) 型时。

这由推论 1 与第六章图版 I 至 IX 得出。

注. 1) 设 \(e_1 = (\mathfrak{g}, \mathfrak{h}_1, \mathrm{B}_1, (X_\alpha^1)_{\alpha \in \mathrm{B}_1})\) ，\(e_2 = (\mathfrak{g}, \mathfrak{h}_2, \mathrm{B}_2, (X_\alpha^2)_{\alpha \in \mathrm{B}_2})\) ，\(e_2' = (\mathfrak{g}, \mathfrak{h}_2, \mathrm{B}_2, (Y_\alpha^2)_{\alpha \in \mathrm{B}_2})\) 是 \(\mathfrak{g}\) 的标架，而 \(s\) （相应地 \(s'\) ）是 \(\operatorname{Aut}_0(\mathfrak{g})\) 中把 \(e_1\) 变为 \(e_2\) （相应地 \(e_2'\) ）的元素。则 \(s|_{\mathfrak{h}_1} = s'|_{\mathfrak{h}_1}\) 。事实上，\(s'^{-1}s \in \operatorname{Aut}_0(\mathfrak{g}, \mathfrak{h}_1)\) 且 \(s'^{-1}s(\mathrm{B}_1) = \mathrm{B}_1\) ，故 \(\varepsilon(s'^{-1}s) = 1\) 。

\begin{enumerate}
\def\labelenumi{\arabic{enumi})}
\setcounter{enumi}{1}
\tightlist
\item
  设 X 是偶 \((\mathfrak{h},\mathrm{B})\) 组成的集合，其中 h 是 g 的分裂嘉当子代数，B 是 \((\mathfrak{g},\mathfrak{h})\) 的根系的一个基。若 \(x=(\mathfrak{h},\mathrm{B})\) 与 \(x'=(\mathfrak{h}',\mathrm{B}')\) 是 X 的两个元素，则存在 \(s\in\operatorname{Aut}_{0}(\mathfrak{g})\) 把 x 变为 \(x'\) （命题 5），且 s 在 h 上的限制 \(s_{x',x}\) 与 s 的选取无关（注 1）。特别地，\(s_{x'',x'}\circ s_{x',x}=s_{x'',x}\) 若 \(x,x',x''\in X\) ，且 \(s_{x,x}=1\) 。满足条件的诸族 \((h_{x})_{x\in\mathrm{X}}\) 组成的集合
\end{enumerate}

\[
a) h _ {x} \in \mathfrak {h} \text {   if   } x = (\mathfrak {h}, B)
\]

\[
b) s _ {x ^ {\prime}, x} (h _ {x}) = h _ {x ^ {\prime}} \text { if } x, x ^ {\prime} \in X
\]

该集合以自然的方式成为一个向量空间 \(\mathfrak{h}(\mathfrak{g})\) ，我们有时称它为 g 的典则嘉当子代数。对 \(x = (\mathfrak{h}, \mathrm{B})\) 与 \(x' = (\mathfrak{h}', \mathrm{B}')\) ，\(s_{x', x}\) 把 B 变为 \(B'\) ，从而把 \((\mathfrak{g}, \mathfrak{h})\) 的根系变为 \((\mathfrak{g}, \mathfrak{h}')\) 的根系；由此，\(\mathfrak{h}(\mathfrak{g})^*\) （\(\mathfrak{h}(\mathfrak{g})\) 的对偶）自然地配备了一个根系 \(\mathrm{R}(\mathfrak{g})\) 以及一个基 \(\mathrm{B}(\mathfrak{g})\) （\(\mathrm{R}(\mathfrak{g})\) 的基）。我们有时称 \(\mathrm{R}(\mathfrak{g})\) 为 g 的典则根系，称 \(\mathrm{B}(\mathfrak{g})\) 为它的典则基。群 \(\mathrm{Aut}(\mathfrak{g})\) 作用在 \(\mathfrak{h}(\mathfrak{g})\) 上，使 \(\mathrm{R}(\mathfrak{g})\) 与 \(\mathrm{B}(\mathfrak{g})\) 保持稳定；\(\mathrm{Aut}(\mathfrak{g})\) 中平凡地作用在 \(\mathfrak{h}(\mathfrak{g})\) 上的元素正是 \(\mathrm{Aut}_0(\mathfrak{g})\) 的元素。

命题 6. 设 \(\mathfrak{h}\) 是 \(\mathfrak{g}\) 的一个分裂嘉当子代数。以第 1 节的记号，我们有 \(\mathrm{Aut}_0(\mathfrak{g}) = \mathrm{Aut}_e(\mathfrak{g}).\mathrm{Ker}\varepsilon = \mathrm{Aut}_e(\mathfrak{g}).f(\mathrm{T}_Q)\) 。

由 §3 第 3 节 命题 10 的推论，\(\mathrm{Aut}_0(\mathfrak{g}) = \mathrm{Aut}_e(\mathfrak{g}).\mathrm{Aut}_0(\mathfrak{g},\mathfrak{h})\) 。另一方面，由 §2 第 2 节 定理 2 的推论，\(\varepsilon (\mathrm{Aut}_e(\mathfrak{g},\mathfrak{h}))\supset \mathrm{W}(\mathrm{R})\) ，故 \(\mathrm{Aut}_0(\mathfrak{g},\mathfrak{h}) = \mathrm{Aut}_e(\mathfrak{g},\mathfrak{h}).\mathrm{Ker}\varepsilon\) 。

注 3. 命题 6 表明，典则同态

\[
\iota : \mathrm{T} _ {\mathrm{Q}} / \operatorname{Im} (\mathrm{T} _ {\mathrm{P}}) \to \operatorname{Aut} _ {0} (\mathfrak {g}) / \operatorname{Aut} _ {e} (\mathfrak {g}),
\]

（由第 2 节的图诱导的）是满射。特别地，\(\mathrm{Aut}_e(\mathfrak{g})\) 包含 \(\mathrm{Aut}_0(\mathfrak{g})\) 的换位子群；我们将看到（§11 第 2 节 命题 3）二者实际上相等。此外，可以证明 \(\iota\) 是单射，换言之

\[
f (\mathrm{T} _ {\mathrm{Q}}) \cap \operatorname{Aut} _ {e} (\mathfrak {g}) = f ^ {\prime} (\mathrm{T} _ {\mathrm{P}}),
\]

（参见 §7 习题 26 d)）。

命题 7. 设 g 是一个可分裂半单李代数，b 是 g 的一个 Borel 子代数，\(p_{1}\) 与 \(p_{2}\) 是包含 b 的两个不同的抛物子代数。则 \(p_{1}\) 与 \(p_{2}\) 在 \(\operatorname{Aut}_{0}(\mathfrak{g})\) 下不共轭。

我们可以假设 \(k\) 是代数闭的。设 \(s \in \mathrm{Aut}_0(\mathfrak{g})\) 满足 \(s(\mathfrak{p}_1) = \mathfrak{p}_2\) 。设 \(\mathfrak{h}\) 是 \(\mathfrak{g}\) 的一个包含在 \(\mathfrak{b} \cap s(\mathfrak{b})\) 中的嘉当子代数（§3 第 3 节 命题 10）。由于 \(\mathfrak{h}\) 与 \(s(\mathfrak{h})\) 都是 \(s(\mathfrak{b})\) 的嘉当子代数，存在 \(u \in [\mathfrak{b}, \mathfrak{b}]\) 使 \(e^{\mathrm{ad}u}(\mathfrak{h}) = s(\mathfrak{h})\) （第七章 §3 第 4 节 定理 3）。把 \(s\) 换成 \(e^{-\mathrm{ad}u}s\) ，我们化归到 \(s(\mathfrak{h}) = \mathfrak{h}\) 的情形，于是 \(s\) 在 \(\mathfrak{h}\) 上诱导出 Weyl 群 W 的一个元素 \(\sigma\) （W 是 \((\mathfrak{g}, \mathfrak{h})\) 的 Weyl 群）（命题 4）。设 C 是对应于 \(\mathfrak{b}\) 的 Weyl 房。则 \(\mathfrak{p}_1\) 与 \(\mathfrak{p}_2\) 对应于面 \(F_1\) 与 \(F_2\) （\(\mathfrak{h}_{\mathbf{R}}\) 的包含在 C 的闭包中的面）。我们有 \(\sigma(F_1) = F_2\) 。由于 \(\sigma \in W\) ，这蕴含 \(F_1 = F_2\) （第五章 §3 第 3 节 定理 2），故 \(\mathfrak{p}_1 = \mathfrak{p}_2\) 。

注 4. 设 g 是一个可分裂半单李代数，P 是 g 的抛物子代数组成的集合，\(\operatorname{Aut}_{0}(\mathfrak{g})\) 作用在这个集合上。沿用注 2 的记号。设 \(\Sigma\) 是 \(\mathrm{B}(\mathfrak{g})\) 的一个子集。给定 \(\Sigma\) 等价于对每个 \(x = (\mathfrak{h}, \mathrm{B}) \in \mathrm{X}\) 给定 B 的一个子集 \(\Sigma_{x}\) ，使 \(s_{x',x}\) 把 \(\Sigma_{x}\) 变为 \(\Sigma_{x'}\) 对任何 \(x, x' \in X\) 成立。设 \(p_{x}\) 是对应于 \(\Sigma_{x}\) 的 g 的抛物子代数（§3 第 4 节 注）。\(p_{x}\) 在 \(\operatorname{Aut}_{0}(\mathfrak{g})\) 作用下的轨道是诸 \(p_{x'}\) （\(x' \in X\) ）组成的集合。这就定义了一个从 \(\mathfrak{P}(\mathrm{B}(\mathfrak{g}))\) 到 \(\mathcal{P}/\operatorname{Aut}_{0}(\mathfrak{g})\) 的映射。这个映射由 §3 第 4 节的注知是满射，由命题 7 知是单射。

\subsection{Aut(g) 上的 Zariski 拓扑}

命题 8. 设 V 是向量空间 g 的自同态组成的集合。则 \(\operatorname{Aut}(\mathfrak{g})\) 在 V 中对 Zariski 拓扑是闭的（第七章附录 I）。

设 \(K\) 是 \(\mathfrak{g}\) 的 Killing 型。若 \(s \in \operatorname{Aut}(\mathfrak{g})\) ，

\[
  [ s x, s y ] = [ x, y ]\tag{2}
\]

\[
\mathrm{K} (s x, s y) = \mathrm{K} (x, y)\tag{3}
\]

对一切 \(x, y \in g\) 成立。反之，设 s 是 V 的元素，满足对一切 \(x, y \in g\) 的 (2) 与 (3)。则 \(\operatorname{Ker}(s) = 0\) ，故 s 是双射且 \(s \in \operatorname{Aut}(\mathfrak{g})\) 。但是，对一切 \(x, y \in g\) ，从 V 到 g 与 k 的映射 \(s \mapsto [sx, sy]\) 与 \(s \mapsto \operatorname{K}(sx, sy)\) 是多项式的。

命题 9. 设 \(\mathfrak{h}\) 是 \(\mathfrak{g}\) 的一个分裂嘉当子代数。

\begin{enumerate}
\def\labelenumi{(\roman{enumi})}
\item
  群 \(f(\mathrm{T}_{\mathrm{Q}})\) 在 \(\operatorname{Aut}(\mathfrak{g})\) 中对 Zariski 拓扑是闭的。
\item
  群 \(f(q(\mathrm{T_P}))\) 在 \(f(\mathrm{T_Q})\) 中对 Zariski 拓扑是稠密的。
\end{enumerate}

断言 (i) 由等式 \(f(\mathrm{T}_{\mathrm{Q}}) = \operatorname{Aut}(\mathfrak{g},\mathfrak{h})\cap \operatorname{Ker}\varepsilon\) （命题 2）得出。置 \(m = (\mathrm{P}(\mathrm{R}):\mathrm{Q}(\mathrm{R}))\) 。设 \(\mathrm{F}\) 是 \(\mathrm{V}\) 上的一个多项式函数；我们假设 \(\mathrm{F}\) 在每个元素的 \(m\) 次幂（元素属于 \(f(\mathrm{T}_{\mathrm{Q}})\) ）上为零，并证明 \(\mathrm{F}|f(\mathrm{T}_{\mathrm{Q}}) = 0\) ；鉴于引理 3，这就将证明 (ii)。

集合 \(V'\) （V 中在 h 上诱导恒等且使每个 \(g^{\alpha}\) 稳定的元素组成的集合）可以与 \(k^{R}\) 等同。设 \(F'\) 是 F 在 \(V' = k^{R}\) 上的限制；这是一个多项式函数。我们有 \(f(\mathrm{T}_{\mathrm{Q}}) \subset \mathrm{V}'\) 。设 \(\mathrm{B} = (\alpha_{1}, \ldots, \alpha_{l})\) 是 R 的一个基。对一切 \(t = (t_{1}, \ldots, t_{l}) \in k^{*B}\) ，用 \(\varphi(t)\) 表示从 \(\mathrm{Q}(\mathrm{R})\) 到群 \(k^{*}\) 的扩张 t 的同态。则 \(\mathrm{F}'(f(\varphi(t)))\) 可以写成有限和

\[
\sum_ {n _ {1}, \dots , n _ {l} \in \mathbf {Z}} c _ {n _ {1}, \dots , n _ {l}} t _ {1} ^ {n _ {1}} \dots t _ {l} ^ {n _ {l}} = \mathrm{H} (t _ {1}, \dots , t _ {l}).
\]

由假设，

\[
0 = \mathrm{H} (t _ {1} ^ {m}, \dots , t _ {l} ^ {m}) = \sum_ {n _ {1}, \dots , n _ {l} \in \mathbf {Z}} c _ {n _ {1}, \dots , n _ {l}} t _ {1} ^ {m n _ {1}} \dots t _ {l} ^ {m n _ {l}}
\]

对一切 \(t_{1},\ldots,t_{l}\in k^{*}\) 成立。于是诸 \(c_{n_{1},\ldots,n_{l}}\) 是一个在 \(k^{*l}\) 上为零的 l 个变量的多项式的系数；因此它们全为零。

命题 10. 假设 \(\mathfrak{g}\) 是可分裂的。

\begin{enumerate}
\def\labelenumi{(\roman{enumi})}
\item
  群 \(\operatorname{Aut}_{e}(\mathfrak{g})\) 在 \(\operatorname{Aut}_{0}(\mathfrak{g})\) 中对 Zariski 拓扑是稠密的。
\item
  群 \(\mathrm{Aut}_e(\mathfrak{g})\) 与 \(\mathrm{Aut}_0(\mathfrak{g})\) 对 Zariski 拓扑是连通的。
\end{enumerate}

由命题 3，\(f(q(\mathrm{T}_{\mathrm{P}})) \subset \mathrm{Aut}_e(\mathfrak{g})\) 。对一切 \(s \in \mathrm{Aut}_e(\mathfrak{g})\) ，\(s.f(q(\mathrm{T}_{\mathrm{P}}))\) 在 Zariski 拓扑中的闭包由命题 9 知包含 \(s.f(\mathrm{T}_{\mathrm{Q}})\) 。因此 \(\mathrm{Aut}_e(\mathfrak{g})\) 的闭包包含 \(\mathrm{Aut}_e(\mathfrak{g}).f(\mathrm{T}_{\mathrm{Q}}) = \mathrm{Aut}_0(\mathfrak{g})\) （命题 6）。这就证明了 (i)。

设 \(\mathrm{Aut}_e(\mathfrak{g}) = \Omega \cup \Omega'\) 是 \(\mathrm{Aut}_e(\mathfrak{g})\) 的一个分割，由 Zariski 拓扑中的相对开子集组成，且 \(\Omega \neq \emptyset\) 。若 \(\omega \in \Omega\) 且 \(x\) 是 \(\mathfrak{g}\) 的幂零元素，则映射 \(\tau : t \mapsto \omega \exp(t \operatorname{ad} x)\) （从 \(k\) 到 \(\mathrm{Aut}_e(\mathfrak{g})\) 的）是多项式的，故在 Zariski 拓扑中连续；于是 \(\tau(k)\) 是连通的；由于 \(\omega \in \tau(k)\) ，我们有 \(\tau(k) \subset \Omega\) 。这样，\(\Omega. (\exp \operatorname{ad} kx) \subset \Omega\) ，故 \(\Omega. \mathrm{Aut}_e(\mathfrak{g}) \subset \Omega\) 且 \(\Omega = \mathrm{Aut}_e(\mathfrak{g})\) 。这就证明 \(\mathrm{Aut}_e(\mathfrak{g})\) 是连通的。于是由 (i)，\(\mathrm{Aut}_0(\mathfrak{g})\) 是连通的。证毕。

我们将看到（§8 第 4 节 命题 6 的推论），\(\mathrm{Aut}_0(\mathfrak{g})\) 在 V 中对 Zariski 拓扑是闭的，且它是 \(\mathrm{Aut}(\mathfrak{g})\) 的单位元的连通分支。另一方面，\(\mathrm{Aut}_e(\mathfrak{g})\) 一般说来在 Zariski 拓扑中不是闭的。

* 假设 \((\mathfrak{g},\mathfrak{h})\) 是分裂的。群 \(\mathrm{Aut}_0(\mathfrak{g})\) 是群 \(G(k)\) ，即 \(k\) 上点的群，其中 \(G\) 是一个具有平凡中心的连通半单代数群（伴随群）。群 \(f(T_{Q})\) 等于 \(H(k)\) ，其中 \(H\) 是 \(G\) 的以 \(\mathfrak{h}\) 为李代数的嘉当子群。逆像 \(\tilde{H}\) （\(H\) 在泛覆盖 \(\tilde{G}\) （\(G\) 的，代数意义上的）中的逆像）以 \(T_{P}\) 为其 \(k\) -点群。\(\tilde{G}(k)\) 在 \(G(k) = \mathrm{Aut}_0(\mathfrak{g})\) 中的像是群 \(\mathrm{Aut}_e(\mathfrak{g})\) 。*

\subsection{李群情形}

命题 11. 假设 k 是 R 、C 或一个非离散完备超度量域。设 h 是 g 的一个分裂嘉当子代数。

\begin{enumerate}
\def\labelenumi{(\roman{enumi})}
\item
  \(\operatorname{Aut}(\mathfrak{g},\mathfrak{h})\) 是 \(\operatorname{Aut}(\mathfrak{g})\) 的一个李子群，其李代数是 ad h。
\item
  \(f(\mathrm{T}_{\mathrm{Q}})\) 与 \((q\circ f)(\mathrm{T}_{\mathrm{P}})\) 是 \(\operatorname {Aut}(\mathfrak{g},\mathfrak{h})\) 的开子群
\item
  \(\mathrm{Aut}_e(\mathfrak{g})\) 是 \(\mathrm{Aut}(\mathfrak{g})\) 的一个开子群。
\item
  若 \(k = \mathbf{R}\) 或 \(\mathbf{C}\) ，则 \(\operatorname{Aut}_e(\mathfrak{g})\) 是 \(\operatorname{Aut}(\mathfrak{g})\) 的单位元连通分支，换言之即 \(\operatorname{Int}(\mathfrak{g})\) 。
\end{enumerate}

由第三章 §3 第 8 节 命题 29 的推论 2 与第 10 节 命题 36，\(\operatorname{Aut}(\mathfrak{g}, \mathfrak{h})\) 是 \(\operatorname{Aut}(\mathfrak{g})\) 的李子群，其李代数是这样的 \(\operatorname{ad} x (x \in \mathfrak{g})\) 组成的集合：使 \((\operatorname{ad} x) \mathfrak{h} \subset \mathfrak{h}\) ，换言之即 \(\operatorname{ad} \mathfrak{h}\) 。

设 \(\mathrm{H} \in \mathfrak{h}\) 。存在 \(\varepsilon > 0\) 具有下列性质：对 \(t \in k\) 与 \(|t| < \varepsilon\) ，\(\exp(t\gamma(\mathrm{H}))\) 对一切 \(\gamma \in \mathrm{P}(\mathrm{R})\) 有定义，且映射 \(\gamma \mapsto \exp(t\gamma(\mathrm{H}))\) 是同态 \(\sigma_t\) （从 \(\mathrm{P}(\mathrm{R})\) 到 \(k^*\) 的）。对 \(|t| < \varepsilon\) ，\(\exp(t\operatorname{ad}\mathrm{H})\) 有定义，在 \(\mathfrak{h}\) 上诱导恒等，且在 \(\mathfrak{g}^\alpha\) 上诱导比为 \(\sigma_t(\alpha)\) 的位似；故 \(\exp t\operatorname{ad}\mathrm{H} \in (q \circ f)(\mathrm{T}_{\mathrm{P}})\) 。这就证明，鉴于 (i)，\((q \circ f)(\mathrm{T}_{\mathrm{P}})\) 包含 \(\operatorname{Aut}(\mathfrak{g},\mathfrak{h})\) 中 1 的一个邻域，从而是 \(\operatorname{Aut}(\mathfrak{g},\mathfrak{h})\) 的开子群。更有甚者，\(f(\mathrm{T}_{\mathrm{Q}})\) 是 \(\operatorname{Aut}(\mathfrak{g},\mathfrak{h})\) 的开子群。

对一切 \(\alpha \in \mathrm{R}\) ，\(\exp \operatorname{ad} \mathfrak{g}^{\alpha} \subset \operatorname{Aut}_{e}(\mathfrak{g})\) 。鉴于 (ii)，\(\operatorname{Aut}_{e}(\mathfrak{g})\) 包含 \(\operatorname{Aut}(\mathfrak{g})\) 中 1 的一个邻域，这就证明了 (iii)。

假设 \(k = \mathbf{R}\) 或 \(\mathbf{C}\) 。则 \(\mathrm{Aut}_e(\mathfrak{g})\) 包含在单位元连通分支 \(\mathrm{C}\) （\(\mathrm{Aut}(\mathfrak{g})\) 的）中（第七章 §3 第 1 节），并由 (iii) 知在 \(\mathrm{Aut}(\mathfrak{g})\) 中是开的。于是 \(\mathrm{Aut}_e(\mathfrak{g}) = \mathrm{C}\) 。最后，由第三章 §9 第 8 节 命题 30 (i)，\(\mathrm{C} = \mathrm{Int}(\mathfrak{g})\) 。

\section{分裂半单李代数上的模}

在本节中，\((\mathfrak{g},\mathfrak{h})\) 表示一个分裂半单李代数，R 表示它的根系，W 表示它的 Weyl 群，B 表示 R 的一个基，\(R_{+}\) （相应地 \(R_{-}\) ）表示相对于 B 的正根（相应地负根）组成的集合。置

\[
\mathfrak {n} _ {+} = \sum_ {\alpha \in R _ {+}} \mathfrak {g} ^ {\alpha}, \quad \mathfrak {n} _ {-} = \sum_ {\alpha \in R _ {-}} \mathfrak {g} ^ {\alpha}, \quad \mathfrak {b} _ {+} = \mathfrak {h} + \mathfrak {n} _ {+} \text {   and   } \mathfrak {b} _ {-} = \mathfrak {h} + \mathfrak {n} _ {-}.
\]

我们有 \(\mathfrak{n}_{+} = [\mathfrak{b}_{+},\mathfrak{b}_{+}],\mathfrak{n}_{-} = [\mathfrak{b}_{-},\mathfrak{b}_{-}].\)

对一切 \(\alpha\in\mathrm{R}\) ，选取一个元素 \(X_{\alpha}\in\mathfrak{g}^{\alpha}\) ，使

\[
[ X _ {\alpha}, X _ {- \alpha} ] = - H _ {\alpha}
\]

（§2 第 4 节）；下面的定义都不依赖这一选取。

\subsection{权与本原元}

设 V 是一个 g-模。对一切 \(\lambda\in h^{*}\) ，用 \(V^{\lambda}\) 表示把 V 视为 h-模时相对于 \(\lambda\) 的准素子空间（第七章 §1 第 1 节）。\(V^{\lambda}\) 的元素称为 g-模 V 的权为 \(\lambda\) 的元素。诸 \(V^{\lambda}\) 的和是直和（第七章 §1 第 1 节 命题 3）。对一切 \(\alpha\in h^{*}\) 与 \(\lambda\in h^{*}\) ，\(g^{\alpha}V^{\lambda}\subset V^{\alpha+\lambda}\) （第七章 §1 第 3 节 命题 10 (ii)）。\(V^{\lambda}\) 的维数称为 \(\lambda\) 在 V 中的重数；若它 \(\geq1\) ，即若 \(V^{\lambda}\neq0\) ，则称 \(\lambda\) 是 V 的一个权。若 V 是有限维的，由 h 的元素定义的 V 的位似是半单的，故 \(V^{\lambda}\) 是由这样的 \(x\in V\) 组成的集合：\(Hx=\lambda(H)x\) 对一切 \(H\in h\) 成立。

引理 1. 设 \(\mathrm{V}\) 是一个 \(\mathfrak{g}\) -模，\(v \in \mathrm{V}\) 。下列条件等价：

\begin{enumerate}
\def\labelenumi{(\roman{enumi})}
\item
  \(\mathfrak{b}_{+}v\subset kv;\)
\item
  \(\mathfrak{h}v \subset kv\) 且 \(n_{+}v = 0;\)
\item
  \(\mathfrak{h}v\subset kv\) 且 \(\mathfrak{g}^{\alpha}v = 0\) 对一切 \(\alpha \in \mathrm{B}\) 成立。
\item
  \(\Longrightarrow\) (ii): 假设 \(\mathfrak{b}_{+}v\subset kv\) 。存在 \(\lambda \in \mathfrak{h}^*\) 使 \(v\in V^{\lambda}\) 。设 \(\alpha \in \mathrm{R}_{+}\) 。则 \(\mathfrak{g}^{\alpha}.v\subset V^{\lambda}\cap V^{\lambda +\alpha} = 0\) 。故 \(\mathfrak{n}_{+}v = 0\) 。
\item
  \(\Longrightarrow\) (iii): 这是显然的。
\item
  \(\Longrightarrow\) (i): 这由 \((X_{\alpha})_{\alpha \in \mathrm{B}}\) 生成 \(\mathfrak{n}_{+}\) 这一事实得出（§3 第 3 节 命题 9 (iii)）。
\end{enumerate}

定义 1. 设 V 是一个 g-模，\(v \in V\) 。若 \(v \neq 0\) 且 v 满足引理 1 的条件，则称 v 是 V 的一个本原元。

本原元属于诸 \(V^{\lambda}\) 之一。对一切 \(\lambda\in h^{*}\) ，用 \(V_{\pi}^{\lambda}\) 表示由这样的 \(v\in V^{\lambda}\) 组成的集合：\(b_{+}v\subset kv\) 。于是，权为 \(\lambda\) 的本原元就是 \(V_{\pi}^{\lambda}\) 的非零元素。

命题 1. 设 V 是一个 g-模，v 是 V 的一个本原元，\(\omega\) 是 v 的权。假设 V 作为 g-模由 v 生成。

\begin{enumerate}
\def\labelenumi{(\roman{enumi})}
\item
  若用 \(\mathrm{U}(\mathfrak{n}_{-})\) 表示 \(\mathfrak{n}_{-}\) 的包络代数，则我们有 \(\mathrm{V} = \mathrm{U}(\mathfrak{n}_{-}).v\) 。
\item
  对一切 \(\lambda \in \mathfrak{h}^*\) ，\(V^{\lambda}\) 是由这样的 \(x \in V\) 组成的集合：\(Hx = \lambda(H)x\) 对一切 \(H \in \mathfrak{h}\) 成立。我们有 \(V = \bigoplus_{\lambda \in \mathfrak{h}^*} V^{\lambda}\) ，且每个 \(V^{\lambda}\) 是有限维的。空间 \(V^{\omega}\) 的维数是 1 ，且 \(V\) 的每个权都形如 \(\omega - \sum_{\alpha \in B} n_{\alpha}. \alpha\) ，其中诸 \(n_{\alpha}\) 是 \(\geq 0\) 的整数。
\item
  V 是一个不可分解的 \(\mathfrak{g}\) -模，且它的交换子约化为纯量。
\item
  设 \(\mathrm{U}(\mathfrak{g})\) 是 g 的包络代数，Z 是 \(\mathrm{U}(\mathfrak{g})\) 的中心。存在唯一的从 Z 到 k 的同态 \(\chi\) ，使对一切 \(z \in Z\) ，\(z_{V}\) 是比为 \(\chi(z)\) 的位似。
\end{enumerate}

设 \(\mathrm{U}(\mathfrak{b}_{+})\) 是 \(\mathfrak{b}_{+}\) 的包络代数。我们有 \(\mathrm{U}(\mathfrak{g}) = \mathrm{U}(\mathfrak{n}_{-})\) . \(\mathrm{U}(\mathfrak{b}_{+})\) （第一章 §2 第 7 节 定理 1 的推论 6）。于是

\[
\mathrm{V} = \mathrm{U} (\mathfrak {g}). v = \mathrm{U} (\mathfrak {n} _ {-}). \mathrm{U} (\mathfrak {b} _ {+}). v = \mathrm{U} (\mathfrak {n} _ {-}). v.
\]

用 \(\alpha_{1},\ldots,\alpha_{n}\) 表示 \(R_{+}\) 的不同元素。则

\[
\left(X _ {- \alpha_ {1}} ^ {p _ {1}} X _ {- \alpha_ {2}} ^ {p _ {2}} \dots X _ {- \alpha_ {n}} ^ {p _ {n}}\right) (p _ {1}, \ldots , p _ {n}) \in \mathbf {N} ^ {n}
\]

是 \(\mathrm{U}(\mathfrak{n}_{-})\) 的一个基，故

\[
\mathrm{V} = \sum_ {(p _ {1}, \dots , p _ {n}) \in \mathbf {N} ^ {n}} k X _ {- \alpha_ {1}} ^ {p _ {1}} \dots X _ {- \alpha_ {n}} ^ {p _ {n}} v.\tag{1}
\]

对 \(\lambda \in \mathfrak{h}^*\) ，置

\[
\mathrm{T} _ {\lambda} = \sum_ {(p _ {1}, \ldots , p _ {n}) \in \mathbf {N} ^ {n}, \omega - p _ {1} \alpha_ {1} - \dots - p _ {n} \alpha_ {n} = \lambda} k X _ {- \alpha_ {1}} ^ {p _ {1}} \ldots X _ {- \alpha_ {n}} ^ {p _ {n}} v.
\]

由第七章 §1 第 1 节 命题 2 (ii)，若 \(h \in \mathfrak{h}\) ，则 \(h_{\mathrm{V}}|_{\mathrm{T}_{\lambda}}\) 是比为 \(\lambda(h)\) 的位似。故 \(\mathrm{T}_{\lambda} \subset \mathrm{V}^{\lambda}\) 。另一方面，(1) 蕴含

\[
\mathrm{V} = \sum_ {\lambda \in \omega - \mathbf {N} \alpha_ {1} - \dots - \mathbf {N} \alpha_ {n}} \mathrm{T} _ {\lambda}.
\]

诸 \(V^{\lambda}\) 的和是直和（第七章 §1 第 1 节 命题 3）。由这些观察得出：\(V^{\lambda} = T_{\lambda}\) ；V 是诸 \(V^{\lambda}\) 的直和；且 \(V^{\lambda}\) 是由这样的 \(x \in V\) 组成的集合：\(hx = \lambda(h)x\) 对一切 \(h \in h\) 成立。另一方面，\(\dim V^{\lambda}\) 至多是这样的 \((p_{1}, \ldots, p_{n}) \in \mathbf{N}^{n}\) 组成的集合的基数：它们满足 \(p_{1}\alpha_{1} + \cdots + p_{n}\alpha_{n} = \omega - \lambda\) 。这就证明：\(V^{\lambda} = 0\) （若 \(\omega - \lambda \notin \sum_{\alpha \in B} N\alpha\) ）；\(\dim V^{\omega} = 1\) ；且诸 \(V^{\lambda}\) 都是有限维的。

设 \(c\) 是 V 的交换子的一个元素。对一切 \(h \in \mathfrak{h}\) ，

\[
h c (v) = c h (v) = \omega (h) c (v),
\]

故 \(c(v) \in V^{\omega}\) ；于是存在 \(t \in k\) 使 \(c(v) = tv\) 。现在，对一切 \((p_1, \ldots, p_n) \in \mathbf{N}^n\) ，

\[
c X _ {- \alpha_ {1}} ^ {p _ {1}} \dots X _ {- \alpha_ {n}} ^ {p _ {n}} v = X _ {- \alpha_ {1}} ^ {p _ {1}} \dots X _ {- \alpha_ {n}} ^ {p _ {n}} c v = t X _ {- \alpha_ {1}} ^ {p _ {1}} \dots X _ {- \alpha_ {n}} ^ {p _ {n}} v
\]

故 \(c = t.1\) 。于是，\(V\) 的交换子约化为纯量。这蕴含 (iv) 以及 \(V\) 不可分解这一事实。

定义 2. 命题 1 (iv) 的同态 \(\chi\) 称为 \(\mathfrak{g}\) -模 V 的中心特征标。

命题 2. 设 V 是由权为 \(\omega\) 的本原元 e 生成的 g-模，X 是一个半单 g-模。设 \(\Phi\) 是从 g-模 V 到 g-模 X 的同态组成的集合。则 \(\varphi\mapsto\varphi(e)\) 是从向量空间 \(\Phi\) 到向量空间 \(X_{\pi}^{\omega}\) 的同构。

显然 \(\varphi(e) \in X_{\pi}^{\omega}\) 对一切 \(\varphi \in \Phi\) 成立。若 \(\varphi \in \Phi\) 且 \(\varphi(e) = 0\) ，则 \(\varphi = 0\) ，因为 \(e\) 生成 \(\mathfrak{g}\) -模 V。我们证明：若 \(f\) 是 \(X_{\pi}^{\omega}\) 的非零元素，则存在 \(\varphi \in \Phi\) 使 \(\varphi(e) = f\) 。设 \(X'\) 是 X 的由 \(f\) 生成的子模。由命题 1，\(X'\) 是不可分解的，由于 X 半单，故它是单的。元素 \((e, f)\) 在 \(\mathfrak{g}\) -模 \(V \times X\) 中是本原的。设 N 是 \(V \times X\) 的由 \((e, f)\) 生成的子模。则 \(N \cap X \subset \mathrm{pr}_2(N) = X'\) ，故 \(N \cap X = 0\) 或 \(X'\) ；若 \(N \cap X = X'\) ，则 N 包含线性无关的元素 \((e, f)\) 与 \((0, f)\) ，它们是权为 \(\omega\) 的本原元；这是荒谬的（命题 1），故 \(N \cap X = 0\) 。于是 \(\mathrm{pr}_1|_{N}\) 是从 N 到 V 的单射映射 \(h\) ；这个映射是满射，因为它的像包含 \(e\) 。于是 \(\varphi = \mathrm{pr}_2 \circ h^{-1}\) 是从 \(\mathfrak{g}\) -模 V 到 \(\mathfrak{g}\) -模 X 的、使 \(\varphi(e) = f\) 的同态。

\subsection{具最高权的单模}

回顾固定 B 在 \(h_{Q}^{*}\) 上定义了一个序关系（第六章 §1 第 6 节）。\(h_{Q}^{*}\) 中 \(\geq 0\) 的元素是 B 的元素以 \(\geq 0\) 的有理数为系数的线性组合。

更一般地，我们将考虑元素 \(\lambda, \mu \in \mathfrak{h}^*\) 之间的如下序关系：

\(\lambda - \mu\) 是 B 的元素以 \(\geq 0\) 的有理数为系数的线性组合。

引理 2. 设 V 是一个单 g-模，\(\omega\) 是 V 的一个权。下列条件等价：

\begin{enumerate}
\def\labelenumi{(\roman{enumi})}
\item
  V 的每个权都形如 \(\omega -\mu\) ，其中 \(\mu\) 是一个 \(\geq 0\) 的根权；
\item
  \(\omega\) 是 V 的最高权；
\item
  对一切 \(\alpha \in \mathrm{B}\) ，\(\omega + \alpha\) 不是 V 的权；
\item
  存在权为 \(\omega\) 的本原元。
\item
  \(\Longrightarrow\) (ii) \(\Longrightarrow\) (iii): 这是显然的。
\item
  \(\Longrightarrow\) (iv): 假设条件 (iii) 满足。对一切 \(h\in \mathfrak{h}\) ，
\end{enumerate}

\(\operatorname {Ker}(h_{\mathrm{V}} - \omega (h))\)

是非零的，包含于 \(V^{\omega}\) ，且在 \(h_{V}\) 下稳定。对 dim h 用归纳法，我们看到存在 \(V^{\omega}\) 中的非零 v 使 \(h v \subset kv\) 。条件 (iii) 蕴含 \(n_{+}v = 0\) ，故 v 是本原的。

\begin{enumerate}
\def\labelenumi{(\roman{enumi})}
\setcounter{enumi}{3}
\tightlist
\item
  \(\Longrightarrow\) (i): 设 \(v\) 是权为 \(\omega\) 的本原元。由于 V 是单的，V 由 \(v\) 生成（作为 \(\mathfrak{g}\) -模）。断言 (i) 于是由命题 1 得出。
\end{enumerate}

证毕。

于是，对任何单 g-模，本原元的存在性等价于最高权的存在性，或等价于极大权的存在性。

存在这样的单 \(\mathfrak{sl}(2,\mathbf{C})\) -模 V：对任何嘉当子代数 \(\mathfrak{h}\) （\(\mathfrak{sl}(2,\mathbf{C})\) 的）它都没有权（§1 习题 14 \(f\) )）。这些模在 \(\mathbf{C}\) 上是无限维的（§1 第 3 节 定理 1）。

命题 3. 设 V 是一个单 \(\mathfrak{g}\) -模，具有最高权 \(\omega\) 。

\begin{enumerate}
\def\labelenumi{(\roman{enumi})}
\item
  V 的本原元就是 \(V^{\omega}\) 的非零元素。
\item
  V 作为 \(\mathfrak{h}\) -模是半单的。
\item
  我们有 \(V = \bigoplus_{\lambda \in h^{*}} V^{\lambda}\) 。对一切 \(\lambda \in h^{*}\) ，\(V^{\lambda}\) 是有限维的。我们有 \(\dim V^{\omega} = 1\) 。
\item
  \(\mathfrak{g}\) -模 V 是绝对单的。
\end{enumerate}

断言 (i)、(ii) 与 (iii) 由命题 1 与引理 2 得出。断言 (iv) 由命题 1 (iii) 与代数第八章 §7 第 3 节得出。

推论. 若 \(\mathrm{V}\) 是有限维的，则典则同态 \(\mathrm{U}(\mathfrak{g}) \to \operatorname{End}(\mathrm{V})\) 是满射。

这由 (iv) 得出，参见代数第八章 §3 第 3 节。

命题 4. 设 V 是一个具有最高权 \(\omega\) 的单 g-模，X 是一个半单 g-模，\(X'\) 是 X 中 V 型的同构型分支。则 \(X'\) 是 X 的由 \(X_{\pi}^{\omega}\) 生成的子模。它的长度等于 \(X_{\pi}^{\omega}\) 的维数。

设 \(X''\) 是 X 的由 \(X_{\pi}^{\omega}\) 生成的子模。显然 X 的每个同构于 V 的子模都包含在 \(X''\) 中。故 \(X' \subset X''\) 。另一方面，设 \(\varPhi\) 是从 g-模 V 到 g-模 X 的同态组成的集合。\(X'\) 的长度是 \(\dim_k \varPhi\) （代数第八章 §4 第 4 节），即 \(\dim_k X_{\pi}^{\omega}\) （命题 2）。

\subsection{存在与唯一性定理}

设 \(\lambda \in \mathfrak{h}^*\) 。由于 \(\mathfrak{b}_{+} = \mathfrak{h} \oplus \mathfrak{n}_{+}\) 且 \(\mathfrak{n}_{+} = [\mathfrak{b}_{+}, \mathfrak{b}_{+}]\) ，映射 \(h + n \mapsto \lambda(h)\) （其中 \(h \in \mathfrak{h}\) ，\(n \in \mathfrak{n}_{+}\) ）（从 \(\mathfrak{b}_{+}\) 到 \(k\) 的）是 \(\mathfrak{b}_{+}\) 的一个 1 维表示。用 \(\mathrm{L}_{\lambda}\) 表示 \(k\) -向量空间 \(k\) ，它配备了由这个表示定义的 \(\mathfrak{b}_{+}\) -模结构。设 \(\mathrm{U}(\mathfrak{g})\) ，\(\mathrm{U}(\mathfrak{b}_{+})\) 是 \(\mathfrak{g}\) ，\(\mathfrak{b}_{+}\) 的包络代数，于是 \(\mathrm{U}(\mathfrak{b}_{+})\) 是 \(\mathrm{U}(\mathfrak{g})\) 的一个子代数；回顾 \(\mathrm{U}(\mathfrak{g})\) 是一个自由右 \(\mathrm{U}(\mathfrak{b}_{+})\) -模（第一章 §2 第 7 节 定理 1 的推论 5）。置

\[
\mathrm{Z} (\lambda) = \mathrm{U} (\mathfrak {g}) \otimes_ {\mathrm{U} (\mathfrak {b} _ {+})} \mathrm{L} _ {\lambda}.\tag{2}
\]

则 \(Z(\lambda)\) 是一个左 \(\mathfrak{g}\) -模。用 \(e\) 表示元素 \(1 \otimes 1\) （\(Z(\lambda)\) 中的）。

命题 5. (i) 元素 \(e\) （\(Z(\lambda)\) 的）是权为 \(\lambda\) 的本原元，且生成 \(\mathfrak{g}\) -模 \(Z(\lambda)\) 。

\begin{enumerate}
\def\labelenumi{(\roman{enumi})}
\setcounter{enumi}{1}
\item
  设 \(\mathrm{Z}^{+}(\lambda) = \sum_{\lambda \neq \mu} \mathrm{Z}(\lambda)^{\mu}\) 。\(\mathrm{Z}(\lambda)\) 的每个异于 \(\mathrm{Z}(\lambda)\) 的子模都包含在 \(\mathrm{Z}^{+}(\lambda)\) 中。
\item
  存在最大子模 \(F_{\lambda}\) （\(Z(\lambda)\) 的，异于 \(Z(\lambda)\) 的）。商模 \(Z(\lambda)/F_{\lambda}\) 是单的且具有最高权 \(\lambda\) 。
\end{enumerate}

显然 e 生成 g-模 \(Z(\lambda)\) 。若 \(x \in b_{+}\) ，

\[
x. e = (x. 1) \otimes 1 = (1. x) \otimes 1 = 1 \otimes x. 1 = \lambda (x) (1 \otimes 1) = \lambda (x) e,
\]

由此得 (i)。

\(\mathfrak{h}\) -模 \(Z(\lambda)\) 是半单的（命题 1）。若 \(G\) 是一个 \(\mathfrak{g}\) -子模（\(Z(\lambda)\) 的），则 \(G = \sum_{\mu \in \mathfrak{h}^*}(G \cap Z(\lambda)^\mu)\) 。假设 \(G \cap Z(\lambda)^\lambda \neq 0\) 蕴含 \(G = Z(\lambda)\) ，因为 \(\dim Z(\lambda)^\lambda = 1\) 且 \(e\) 生成 \(\mathfrak{g}\) -模 \(Z(\lambda)\) 。若 \(G \neq Z(\lambda)\) ，则 \(G = \sum_{\mu \neq \lambda} G \cap Z(\lambda)^\mu \subset Z^+( \lambda)\) 。

设 \(F_{\lambda}\) 是 \(Z(\lambda)\) 的异于 \(Z(\lambda)\) 的诸 g-子模之和。由 (ii)，\(F_{\lambda} \subset Z^{+}(\lambda)\) 。故 \(F_{\lambda}\) 是 \(Z(\lambda)\) 的异于 \(Z(\lambda)\) 的最大子模。显然 \(Z(\lambda)/F_{\lambda}\) 是单的，且 e 在 \(Z(\lambda)/F_{\lambda}\) 中的典则像是权为 \(\lambda\) 的本原元。

在本章的其余部分，命题 5 的 g-模 \(Z(\lambda)/F_{\lambda}\) 将记作 \(E(\lambda)\) 。

定理 1. 设 \(\lambda \in \mathfrak{h}^*\) 。\(\mathfrak{g}\) -模 \(\mathrm{E}(\lambda)\) 是单的且具有最高权 \(\lambda\) 。每个单 \(\mathfrak{g}\) -模若其最高权为 \(\lambda\) ，便同构于 \(\mathrm{E}(\lambda)\) 。

第一个断言由命题 5 (iii) 得出。第二个由命题 4 得出。

命题 6. 设 V 是一个 g-模，\(\lambda\) 是 \(h^{*}\) 的一个元素，v 是 V 的权为 \(\lambda\) 的本原元。

\begin{enumerate}
\def\labelenumi{(\roman{enumi})}
\item
  存在唯一的 \(\mathfrak{g}\) -模同态 \(\psi: \mathrm{Z}(\lambda) \to \mathrm{V}\) 使 \(\psi(e) = v\) 。
\item
  假设 v 生成 V。则 \(\psi\) 是满射。此外，\(\psi\) 是双射当且仅当对 \(\mathrm{U}(\mathfrak{n}_{-})\) 的每个非零元素 u ，\(u_{V}\) 是单射。
\item
  映射 \(u \mapsto u \otimes 1\) （从 \(\mathrm{U}(\mathfrak{n}_{-})\) 到 \(\mathrm{Z}(\lambda)\) 的）是双射。
\end{enumerate}

设 K 是 \(\mathrm{U}(\mathfrak{b}_{+})\) 在 \(L_{\lambda}\) 上的表示的核；它在 \(\mathrm{U}(\mathfrak{b}_{+})\) 中的余维数是 1 。设 \(\mathrm{J} = \mathrm{U}(\mathfrak{g})\mathrm{K}\) 是由 K 生成的 \(\mathrm{U}(\mathfrak{g})\) 的左理想；于是 \(L_{\lambda}\) 可以与 \(\mathrm{U}(\mathfrak{b}_{+})/\mathrm{K}\) 等同（作为左 \(\mathrm{U}(\mathfrak{b}_{+})\) -模），且 \(\mathrm{Z}(\lambda)\) 可以与 \(\mathrm{U}(\mathfrak{g})/\mathrm{J}\) 等同（作为左 \(\mathrm{U}(\mathfrak{g})\) -模）。我们有 K.v = 0 ，故 J.v = 0 ，这就证明了 (i)。

现在假设 \(v\) 生成 V。显然 \(\psi\) 是满射。

由 Poincaré-Birkhoff-Witt 定理（第一章 §2 第 7 节 定理 1 的推论 6），\(\mathrm{U}(\mathfrak{n}_{-})\) 在 \(k\) 上的一个基也是 \(\mathrm{U}(\mathfrak{g})\) 作为右 \(\mathrm{U}(\mathfrak{b}_{+})\) -模的一个基。故映射 \(\varphi : u \mapsto u \otimes 1\) （从 \(\mathrm{U}(\mathfrak{n}_{-})\) 到 \(\mathrm{U}(\mathfrak{g}) \otimes_{\mathrm{U}(\mathfrak{b}_{+})} \mathrm{L}_{\lambda}\) 的）是双射。设 \(u \in \mathrm{U}(\mathfrak{n}_{-})\) 。则 \(\varphi^{-1} \circ u_{\mathrm{Z}(\lambda)} \circ \varphi\) 是由 \(u\) 在 \(\mathrm{U}(\mathfrak{n}_{-})\) 上作的左乘。鉴于第一章 §2 第 7 节 定理 1 的推论 7，\(u_{\mathrm{Z}(\lambda)}\) 是单射若 \(u \neq 0\) 。于是，若 \(\psi\) 是双射，则 \(u_{\mathrm{V}}\) 是单射，只要 \(u\) 是 \(\mathrm{U}(\mathfrak{n}_{-})\) 中的非零元。

假设 \(\psi\) 不是单射。存在 \(u \in \mathrm{U}(\mathfrak{n}_{-})\) 使 \(u \neq 0\) 且 \(\psi(\varphi(u)) = 0\) 。则

\[
u _ {\mathrm{V}}. v = u _ {\mathrm{V}}. \psi (1 \otimes 1) = \psi (u \otimes 1) = \psi (\varphi (u)) = 0.
\]

推论 1. 设 \(\lambda \in \mathfrak{h}^*\) 与 \(\alpha \in B\) 使 \(\lambda(H_{\alpha}) + 1 \in \mathbf{N}\) 成立。则 \(Z(-\alpha + s_{\alpha}\lambda)\) 同构于一个 \(\mathfrak{g}\) -子模（\(Z(\lambda)\) 的）。

置 \(m = \lambda(H_{\alpha})\) 。设 \(x = X_{-\alpha}^{m+1}.e \in Z(\lambda)\) ，并设 V 是由 x 生成的 \(Z(\lambda)\) 的子模。则 \(x \neq 0\) （命题 6）。另一方面，\(x \in Z(\lambda)^{\lambda-(m+1)\alpha}\) 。对 \(\beta \in B\) 且 \(\beta \neq \alpha\) ，有 \([\mathfrak{g}^{-\alpha}, \mathfrak{g}^{\beta}] = 0\) 且 \(g^{\beta}.e = 0\) ，故 \(g^{\beta}.x = 0\) 。最后，由于 \([X_{\alpha}, X_{-\alpha}] = -H_{\alpha}\) ，我们有

\[
[ X _ {\alpha}, X _ {- \alpha} ^ {m + 1} ] = (m + 1) X _ {- \alpha} ^ {m} (- H _ {\alpha} + m)
\]

（§1 第 1 节 引理 1 (ii)），故

\[
X _ {\alpha}. x = X _ {\alpha} X _ {- \alpha} ^ {m + 1}. e = [ X _ {\alpha}, X _ {- \alpha} ^ {m + 1} ]. e = (m + 1) X _ {- \alpha} ^ {m} (m e - \lambda (H _ {\alpha}) e) = 0.
\]

于是，\(x\) 是权为 \(\lambda - (m + 1)\alpha\) 的本原元。鉴于命题 6，\(\mathfrak{g}\) -模 \(V\) 同构于 \(Z(-\alpha + \lambda - m\alpha) = Z(-\alpha + s_{\alpha}\lambda)\) 。

推论 2. 设 \(\rho = \frac{1}{2}\sum_{\alpha \in \mathrm{R}_{+}}\alpha\) ，并设 \(\lambda, \mu \in \mathfrak{h}^*\) 。假设 \(\lambda + \rho\) 是 \(\mathrm{R}\) 中的一个支配权，且存在 \(w \in \mathrm{W}\) 使 \(\mu + \rho = w(\lambda + \rho)\) 成立。则 \(\mathrm{Z}(\mu)\) 同构于 \(\mathrm{Z}(\lambda)\) 的一个子模。

当 w = 1 时断言是明显的。假设当 w 的长度 \textless{} q 时断言已经成立。若 w 的长度为 q，则存在 \(\alpha \in B\) 使 \(w = s_{\alpha}w'^{-1}\) 成立，其中 \(l(w') = q - 1\) 。我们有 \(w'(\alpha) \in \mathrm{R}_{+}\) （第六章 §1 第 6 节 命题 17 的推论 2），因而 \(w'^{-1}(\lambda + \rho)(H_{\alpha}) = (\lambda + \rho)(H_{w'\alpha})\) 是一个 \(\geq 0\) 的整数。置

\[
\mu^ {\prime} = w ^ {- 1} (\lambda + \rho) - \rho .
\]

由归纳假设，\(Z(\mu')\) 同构于 \(Z(\lambda)\) 的一个子模。另一方面，由第六章 §1 第 10 节 命题 29 (ii)，

\[
- \alpha + s _ {\alpha} \mu^ {\prime} = - \alpha + s _ {\alpha} w ^ {- 1} (\lambda + \rho) - s _ {\alpha} \rho = w (\lambda + \rho) - \rho = \mu .
\]

此外，\(\rho(H_{\alpha}) = 1\) （第六章 §1 命题 29 (iii)），故 \(\mu'(H_{\alpha}) + 1 \in \mathbf{N}\) 。于是由推论 1 知，\(Z(\mu)\) 同构于 \(Z(\mu')\) 的一个子模，从而也同构于 \(Z(\lambda)\) 的一个子模。

\subsection{\texorpdfstring{\(\mathfrak{h}\) 在 \(\mathfrak{g}\) 的包络代数中的交换子}{\textbackslash mathfrak\{h\} 在 \textbackslash mathfrak\{g\} 的包络代数中的交换子}}

设 U 是 g 的包络代数，\(V \subset U\) 是 h 的包络代数。代数 V 可以与 h 的对称代数 \(\mathbf{S}(\mathfrak{h})\) 等同，也可以与 \(h^{*}\) 上的多项式函数代数等同。用 \(\alpha_{1},\ldots,\alpha_{n}\) 记两两不同的正根。设 \((H_{1},\ldots,H_{l})\) 是 h 的一个基。由 Poincaré-Birkhoff-Witt 定理，诸元素

\[
u ((q _ {i}), (m _ {i}), (p _ {i})) = X _ {- \alpha_ {1}} ^ {q _ {1}} \dots X _ {- \alpha_ {n}} ^ {q _ {n}} H _ {1} ^ {m _ {1}} \dots H _ {l} ^ {m _ {l}} X _ {\alpha_ {1}} ^ {p _ {1}} \dots X _ {\alpha_ {n}} ^ {p _ {n}}
\]

\((q_{i},m_{i},p_{i})\) 为整数 \(\geq 0\) ）构成向量空间 U 的一个基。对一切 \(h\in \mathfrak{h}\) ，我们有

\[
[ h, u ((q _ {i}), (m _ {i}), (p _ {i})) ] = ((p _ {1} - q _ {1}) \alpha_ {1} + \dots + (p _ {n} - q _ {n}) \alpha_ {n}) (h) u ((q _ {i}), (m _ {i}), (p _ {i})).\tag{3}
\]

向量空间 U 在伴随表示下是一个 g-模（从而也是一个 h-模）。若 \(\lambda \in h^{*}\) ，则子空间 \(U^{\lambda}\) 与 \(U_{\lambda}\) 已有定义（第七章 §1 第 3 节）；公式 (3) 表明 \(U^{\lambda} = U_{\lambda}\) ，且 \(U = \bigoplus_{\lambda \in Q} U^{\lambda}\) （其中 Q 是 R 的根权格）。特别地，\(U^{0}\) 是 h 或 V 在 U 中的交换子。

引理 3. 置 \(\mathrm{L} = (\mathfrak{n}_{-}\mathrm{U})\cap \mathrm{U}^{0}\)

\begin{enumerate}
\def\labelenumi{(\roman{enumi})}
\item
  我们有 \(\mathrm{L} = (\mathrm{U}\mathfrak{n}_{+})\cap \mathrm{U}^{0}\) ，且 \(\mathrm{L}\) 是 \(\mathrm{U}^0\) 的一个双边理想。
\item
  我们有 \(U^{0} = V \oplus L\) 。
\end{enumerate}

显然 \(\mathfrak{n}_{-}\mathrm{U}\) （相应地 \(\mathrm{U}\mathfrak{n}_{+}\) ）是这样的元素 \(u((q_{i}),(m_{i}),(p_{i}))\) 的线性组合组成的集合：它们使 \(\sum q_{i}>0\) （相应地 \(\sum p_{i}>0\) ）。另一方面

\[
u \left(\left(q _ {i}\right), \left(m _ {i}\right), \left(p _ {i}\right)\right) \in \mathrm{U} ^ {0} \Longleftrightarrow p _ {1} \alpha_ {1} + \dots + p _ {n} \alpha_ {n} = q _ {1} \alpha_ {1} + \dots + q _ {n} \alpha_ {n}.
\]

这蕴含 \((\mathfrak{n}_{-}\mathrm{U})\cap\mathrm{U}^{0}=(\mathrm{U}\mathfrak{n}_{+})\cap\mathrm{U}^{0}\) 。最后，\((\mathfrak{n}_{-}\mathrm{U})\cap\mathrm{U}^{0}\) （相应地 \((\mathrm{U}\mathfrak{n}_{+})\cap\mathrm{U}^{0}\) ）是 \(U^{0}\) 的一个右（相应地左）理想，由此得 (i)。此外，元素 \(u((q_{i}),(m_{i}),(p_{i}))\) 若属于 \(U^{0}\) ，则它属于 V（相应地属于 L）当且仅当 \(p_{1}=\cdots=p_{n}=q_{1}=\cdots=q_{n}=0\) （相应地 \(p_{1}+\cdots+p_{n}+q_{1}+\cdots+q_{n}>0\) ），由此得 (ii)。证毕。

鉴于引理 3，\(U^{0}\) 到 V 上的、以 L 为核的投影是一个代数同态。它称为从 \(U^{0}\) 到 V 的 Harish-Chandra 同态（相对于 B 的）。回忆 V 可以与 \(h^{*}\) 上的多项式函数代数等同。

命题 7. 设 \(\lambda \in \mathfrak{h}^*\) ，E 是一个 \(\mathfrak{g}\) -模，由权为 \(\lambda\) 的一个本原元生成，\(\chi\) 是 E 的中心特征标，\(\varphi\) 是从 \(\mathrm{U}^0\) 到 V 的 Harish-Chandra 同态。则 \(\chi(z) = (\varphi(z))(\lambda)\) 对 U 的中心中的所有 \(z\) 成立。

设 v 是 E 的权为 \(\lambda\) 的本原元，z 是 U 的中心的一个元素。存在 \(u_{1},\ldots,u_{p}\in U\) 与 \(n_{1},\ldots,n_{p}\in\mathfrak{n}_{+}\) 使 \(z=\varphi(z)+u_{1}n_{1}+\cdots+u_{p}n_{p}\) 成立。则

\[
\chi (z) v = z v = \varphi (z) v + u _ {1} n _ {1} v + \dots + u _ {p} n _ {p} v = \varphi (z) v = (\varphi (z)) (\lambda) v.
\]

推论. 设 \(\langle\cdot,\cdot\rangle\) 是 g 上的一个非退化不变对称双线性型，C 是与 \(\langle\cdot,\cdot\rangle\) 相伴的 Casimir 元。也用 \(\langle\cdot,\cdot\rangle\) 记 \(h^{*}\) 上的一个逆型，即 \(\langle\cdot,\cdot\rangle\) 在 h 上的限制的逆型（§2 第 3 节 命题 5）。则 \(\chi(\mathrm{C})=\langle\lambda,\lambda+2\rho\rangle\) ，其中 \(\rho=\frac{1}{2}\sum_{\alpha\in R_{+}}\alpha\) 。

我们回忆 §2 第 3 节 命题 6 的记号。我们有

\[
\begin{array}{r} \mathrm{C} = \sum_ {\alpha \in \mathrm{R} _ {-}} \frac {1}{\langle X _ {\alpha} , X _ {- \alpha} \rangle} X _ {\alpha} X _ {- \alpha} + \sum_ {\alpha \in \mathrm{R} _ {+}} \frac {1}{\langle X _ {\alpha} , X _ {- \alpha} \rangle} X _ {- \alpha} X _ {\alpha} \\ + \sum_ {\alpha \in \mathrm{R} _ {+}} \frac {1}{\langle X _ {\alpha} , X _ {- \alpha} \rangle} [ X _ {\alpha}, X _ {- \alpha} ] + \sum_ {i \in I} H _ {i} H _ {i} ^ {\prime} \end{array}
\]

故

\[
\varphi (\mathrm{C}) = \sum_ {\alpha \in \mathrm{R} _ {+}} \frac {1}{\langle X _ {\alpha} , X _ {- \alpha} \rangle} [ X _ {\alpha}, X _ {- \alpha} ] + \sum_ {i \in I} H _ {i} H _ {i} ^ {\prime}.
\]

由命题 7，

\[
\chi (\mathrm{C}) = \sum_ {\alpha \in \mathrm{R} _ {+}} \frac {1}{\langle X _ {\alpha} , X _ {- \alpha} \rangle} \lambda ([ X _ {\alpha}, X _ {- \alpha} ]) + \sum_ {i \in I} \lambda (H _ {i}) \lambda (H _ {i} ^ {\prime}).
\]

设 \(h_{\lambda}\) 是 \(\mathfrak{h}\) 中使 \(\langle h_{\lambda}, h \rangle = \lambda(h)\) 对一切 \(h \in \mathfrak{h}\) 成立的元素。由 §2 第 2 节 命题 1，

\[
\lambda \left(\frac {1}{\langle X _ {\alpha} , X _ {- \alpha} \rangle} [ X _ {\alpha}, X _ {- \alpha} ]\right) = \left\langle h _ {\lambda}, \frac {1}{\langle X _ {\alpha} , X _ {- \alpha} \rangle} [ X _ {\alpha}, X _ {- \alpha} ] \right\rangle = \alpha (h _ {\lambda}) = \langle \lambda , \alpha \rangle .
\]

于是

\[
\chi (\mathrm{C}) = \left(\sum_ {\alpha \in \mathrm{R} _ {+}} \langle \lambda , \alpha \rangle\right) + \langle \lambda , \lambda \rangle = \langle \lambda , \lambda + 2 \rho \rangle .
\]

\section{分裂半单李代数上的有限维模}

本节中，我们保留 §6 的一般记号。用 P（相应地 Q）记 R 的权格（相应地根权格）。用 \(P_{+}\) （相应地 \(Q_{+}\) ）记 P（相应地 Q）中对由 B 定义的序关系为正的元素组成的集合。用 \(P_{++}\) 记相对于 B 的 R 的支配权组成的集合（第六章 §1 第 10 节）。\(\lambda\) 是 \(h^{*}\) 的一个元素，它属于 P（相应地属于 \(P_{++}\) ）当且仅当所有的 \(\lambda(H_{\alpha})\) （\(\alpha \in B\) ）都是整数（相应地 \(\geq 0\) 的整数）。我们有 \(P_{++} \subset P_{+}\) （第六章 §1 第 6 节）。若 \(w \in W\) ，用 \(\varepsilon(w)\) 记 w 的行列式，它等于 1 或 -1。置 \(\rho = \frac{1}{2} \sum_{\alpha \in R_{+}} \alpha\) 。